\chapter{环论图片转写待审核稿}

\begin{dxtips}

本文件为图片与 Markdown 笔记整理出的待审核转写稿，已尽量按 elegantbook 环境归类；正式并入正文前仍建议逐条复核。

\end{dxtips}

\section*{7理想}

\begin{note}[来源上级主题]

第三章 (\%E7\%AC\%AC\%E4\%B8\%89\%E7\%AB\%A0\%201ecd2efa3f4280c7a6b0ea0de416c09f.md)

\end{note}

\begin{note}[7理想]

子级 项目: 理想子环定义 (\%E7\%90\%86\%E6\%83\%B3\%E5\%AD\%90\%E7\%8E\%AF\%E5\%AE\%9A\%E4\%B9\%89\%201edd2efa3f4280889eb4daa974cd0880.md), 定理1除环只有两个理想0理想和单位理想 (\%E5\%AE\%9A\%E7\%90\%861\%E9\%99\%A4\%E7\%8E\%AF\%E5\%8F\%AA\%E6\%9C\%89\%E4\%B8\%A4\%E4\%B8\%AA\%E7\%90\%86\%E6\%83\%B30\%E7\%90\%86\%E6\%83\%B3\%E5\%92\%8C\%E5\%8D\%95\%E4\%BD\%8D\%E7\%90\%86\%E6\%83\%B3\%201edd2efa3f4280198c05df2fd987f2b8.md), 主理想 (\%E4\%B8\%BB\%E7\%90\%86\%E6\%83\%B3\%201edd2efa3f42806b9f8cf1209d23b5ab.md), 生成理想符号（a1，a2） (\%E7\%94\%9F\%E6\%88\%90\%E7\%90\%86\%E6\%83\%B3\%E7\%AC\%A6\%E5\%8F\%B7\%EF\%BC\%88a1\%EF\%BC\%8Ca2\%EF\%BC\%89\%201edd2efa3f4280ed86cbd75462c70e23.md)

\end{note}

\begin{note}[来源路径]

近世代数/近世代数/无标题/7理想 1edd2efa3f4280a88bfce32f4c292dfb.md

\end{note}

\section*{8剩余类环同态与理想}

\begin{note}[来源上级主题]

第三章 (\%E7\%AC\%AC\%E4\%B8\%89\%E7\%AB\%A0\%201ecd2efa3f4280c7a6b0ea0de416c09f.md)

\end{note}

\begin{note}[8剩余类环同态与理想]

子级 项目: 定理1 理想和正规子群很像 模理想的剩余类还是环而且和原环同态 (\%E5\%AE\%9A\%E7\%90\%861\%20\%E7\%90\%86\%E6\%83\%B3\%E5\%92\%8C\%E6\%AD\%A3\%E8\%A7\%84\%E5\%AD\%90\%E7\%BE\%A4\%E5\%BE\%88\%E5\%83\%8F\%20\%E6\%A8\%A1\%E7\%90\%86\%E6\%83\%B3\%E7\%9A\%84\%E5\%89\%A9\%E4\%BD\%99\%E7\%B1\%BB\%E8\%BF\%98\%E6\%98\%AF\%E7\%8E\%AF\%E8\%80\%8C\%E4\%B8\%94\%E5\%92\%8C\%E5\%8E\%9F\%E7\%8E\%AF\%E5\%90\%8C\%E6\%80\%81\%201edd2efa3f428071ac13fae2aef057df.md), 定理2 R和自己的商环同态 同态满射核就是理想 (\%E5\%AE\%9A\%E7\%90\%862\%20R\%E5\%92\%8C\%E8\%87\%AA\%E5\%B7\%B1\%E7\%9A\%84\%E5\%95\%86\%E7\%8E\%AF\%E5\%90\%8C\%E6\%80\%81\%20\%E5\%90\%8C\%E6\%80\%81\%E6\%BB\%A1\%E5\%B0\%84\%E6\%A0\%B8\%E5\%B0\%B1\%E6\%98\%AF\%E7\%90\%86\%E6\%83\%B3\%201edd2efa3f4280868a83cd8270b657a6.md), 剩余类环的定义与符号 也就是商环 (\%E5\%89\%A9\%E4\%BD\%99\%E7\%B1\%BB\%E7\%8E\%AF\%E7\%9A\%84\%E5\%AE\%9A\%E4\%B9\%89\%E4\%B8\%8E\%E7\%AC\%A6\%E5\%8F\%B7\%20\%E4\%B9\%9F\%E5\%B0\%B1\%E6\%98\%AF\%E5\%95\%86\%E7\%8E\%AF\%201efd2efa3f428062bb60fef4a28718a3.md), 定理3  R和R的商环同态满射的性质 传递 (\%E5\%AE\%9A\%E7\%90\%863\%20R\%E5\%92\%8CR\%E7\%9A\%84\%E5\%95\%86\%E7\%8E\%AF\%E5\%90\%8C\%E6\%80\%81\%E6\%BB\%A1\%E5\%B0\%84\%E7\%9A\%84\%E6\%80\%A7\%E8\%B4\%A8\%20\%E4\%BC\%A0\%E9\%80\%92\%201edd2efa3f4280a6a5a6cd9145845beb.md)

\end{note}

\begin{note}[来源路径]

近世代数/近世代数/无标题/8剩余类环同态与理想 1edd2efa3f42801295a9f3bc6dd211c2.md

\end{note}

\section*{9极大理想}

\begin{note}[来源上级主题]

第三章 (\%E7\%AC\%AC\%E4\%B8\%89\%E7\%AB\%A0\%201ecd2efa3f4280c7a6b0ea0de416c09f.md)

\end{note}

\begin{note}[9极大理想]

子级 项目: 极大理想的定义 (\%E6\%9E\%81\%E5\%A4\%A7\%E7\%90\%86\%E6\%83\%B3\%E7\%9A\%84\%E5\%AE\%9A\%E4\%B9\%89\%201efd2efa3f4280c6b178ee01ba7ba017.md), 极大理想引理1 对剩余类环充要 (\%E6\%9E\%81\%E5\%A4\%A7\%E7\%90\%86\%E6\%83\%B3\%E5\%BC\%95\%E7\%90\%861\%20\%E5\%AF\%B9\%E5\%89\%A9\%E4\%BD\%99\%E7\%B1\%BB\%E7\%8E\%AF\%E5\%85\%85\%E8\%A6\%81\%201efd2efa3f4280549cb7d1f5b63a5db3.md), 引理2 交换环只有零理想和单位理想R是一个域 (\%E5\%BC\%95\%E7\%90\%862\%20\%E4\%BA\%A4\%E6\%8D\%A2\%E7\%8E\%AF\%E5\%8F\%AA\%E6\%9C\%89\%E9\%9B\%B6\%E7\%90\%86\%E6\%83\%B3\%E5\%92\%8C\%E5\%8D\%95\%E4\%BD\%8D\%E7\%90\%86\%E6\%83\%B3R\%E6\%98\%AF\%E4\%B8\%80\%E4\%B8\%AA\%E5\%9F\%9F\%20202d2efa3f42801f8790d49e461b4508.md), 极大理想与域 (\%E6\%9E\%81\%E5\%A4\%A7\%E7\%90\%86\%E6\%83\%B3\%E4\%B8\%8E\%E5\%9F\%9F\%20205d2efa3f4280fab07ae7a25e907d3b.md)

\end{note}

\begin{note}[来源路径]

近世代数/近世代数/无标题/9极大理想 1efd2efa3f42802fa98fcbcc6a6ed8b2.md

\end{note}

\section*{n阶循环群互相同构 无限循环群彼此同构 命题1.2.9}

\begin{note}[来源上级主题]

1.3有限群 (1\%203\%E6\%9C\%89\%E9\%99\%90\%E7\%BE\%A4\%201ead2efa3f4280c39a49d8b0af9aec03.md)

\end{note}

\begin{proposition}[n阶循环群互相同构 无限循环群彼此同构 命题1.2.9]

Aj 1.26

AG = (x) ZAIRE, GE |x| =n, WG = \{e,x,x72,--- x7 |\}, PPR PA kLAZAAEA

9. AAAS HEY A IRB EY Pe                            “

DETAR AME? WALDL, TEMA PR ORME, ABIL AL O FRSA AT n SRR EMA NTA, FLA nn = |x].

UEW) ATRIA. A, E-SG PRAMAVSRKMAOAMH A MEN BA;s Bo, MOAMN A

ne eA LH

AN RIEHZ-\&. ERG PRK x”, Ht me Z. PREPARE, RNA ma=qntr, HH qeZ

O<r<n-1l, MAAAXx" =e, Px axI =x") x =x", MARAT AO ATH An

AM RIEA BOR. RUE, BHRO<m'<m<n-l, fex™ax™, Mx ™ =e, HEL <m-m'<

n-l<n, Pen=|x| LRDHERKK Hx =e, RRSRT TH.

Se LPR, G=f\{e,x,x7,--- x "\}, SP RB eR EY.

\end{proposition}

\begin{note}[图片 OCR 摘录，待人工复核]
以下内容来自源图片识别，便于审核时对照：

\textbf{图片 1}（image.png）
\begin{Verbatim}[fontsize=\small]
Aj 1.26

AG = (x) ZAIRE, GE |x| =n, WG = {e,x,x72,--- x7 |}, PPR PA kLAZAAEA

9. AAAS HEY A IRB EY Pe                            “

DETAR AME? WALDL, TEMA PR ORME, ABIL AL O FRSA AT n SRR EMA NTA, FLA nn = |x].
UEW) ATRIA. A, E-SG PRAMAVSRKMAOAMH A MEN BA;s Bo, MOAMN A
ne eA LH

AN RIEHZ-&. ERG PRK x”, Ht me Z. PREPARE, RNA ma=qntr, HH qeZ
O<r<n-1l, MAAAXx" =e, Px axI =x") x =x", MARAT AO ATH An

AM RIEA BOR. RUE, BHRO<m'<m<n-l, fex™ax™, Mx ™ =e, HEL <m-m'<
n-l<n, Pen=|x| LRDHERKK Hx =e, RRSRT TH.

Se LPR, G=f{e,x,x7,--- x "}, SP RB eR EY.
\end{Verbatim}

\end{note}

\begin{note}[来源路径]

近世代数/近世代数/无标题/n阶循环群互相同构 无限循环群彼此同构 命题1 2 9 1d2d2efa3f4280ebb6eaf07515b209e7.md

\end{note}

\section*{主理想}

\begin{note}[来源上级主题]

7理想 (7\%E7\%90\%86\%E6\%83\%B3\%201edd2efa3f4280a88bfce32f4c292dfb.md)

\end{note}

\begin{note}[主理想]

主理想要求就是由一个元生成的理想

证明是主理想就是证明一个元素能生成整个环，不是每个理想都能做到的因为理想的诞生是理想中元素与环中元素作用并且不断补充进理想，而生成整个理想是理想里面的一个元素不与环作用

可交换+有单位元=理想可以写成ra形式 r∈环

\end{note}

\begin{note}[关联图片]
以下图片已纳入本条整理，当前版本优先依据 Markdown 文字整理，图片细节可在审核时对照：

1. image.png\\

\end{note}

\begin{note}[来源路径]

近世代数/近世代数/无标题/主理想 1edd2efa3f42806b9f8cf1209d23b5ab.md

\end{note}

\section*{交换环与单位元}

\begin{note}[来源上级主题]

第三章 (\%E7\%AC\%AC\%E4\%B8\%89\%E7\%AB\%A0\%201ecd2efa3f4280c7a6b0ea0de416c09f.md)

\end{note}

\begin{note}[交换环与单位元]

EM —PHR WW-PRRB, BUI R HEB Ta. bR-

aA

ab=ba.

F7E—-TRRAB, MFEREBRRAn UVRAWERA Ta. b Rik,

A

a"b"=(ab)”.

Bi LRCBRNCAAS ST RPMHBRE. EHH NLBRNK

A BR-S RBA PX FHRVIER LY isc. (A—TS ABUT,

RTT UBR, XPSschSs SME BE A Me is.

EM —TPHRWN—PIce WR—-T RMI, SMF R HER K

Bi, BBA

ea=ae\textasciitilde{}=a.

—fxth, —P ARYA PS BIC.

\end{note}

\begin{note}[图片 OCR 摘录，待人工复核]
以下内容来自源图片识别，便于审核时对照：

\textbf{图片 1}（image.png）
\begin{Verbatim}[fontsize=\small]
EM —PHR WW-PRRB, BUI R HEB Ta. bR-
aA
ab=ba.
F7E—-TRRAB, MFEREBRRAn UVRAWERA Ta. b Rik,
A
a"b"=(ab)”.
Bi LRCBRNCAAS ST RPMHBRE. EHH NLBRNK
A BR-S RBA PX FHRVIER LY isc. (A—TS ABUT,
RTT UBR, XPSschSs SME BE A Me is.
EM —TPHRWN—PIce WR—-T RMI, SMF R HER K
Bi, BBA
ea=ae~=a.
—fxth, —P ARYA PS BIC.
\end{Verbatim}

\end{note}

\begin{note}[来源路径]

近世代数/近世代数/无标题/交换环与单位元 1ecd2efa3f4280739601e8d3434ada2a.md

\end{note}

\section*{交换群，循环群与商群}

\begin{note}[交换群，循环群与商群]

5) OBEN E— FAR CB, SURG RE tHe ACRE.

(aN)(bN) =(ab)N =(ba)N =(bN )(aN)

6) PEAR BEHVLE— FAAS EE, LETRA AL BE.

Ck G=(a) X87 N<G

APPA ARRH, ATARI THAN TH, We

NIG.

VbeG,b=a bN=a'N=(aN)’

mh 8 8G/N=\{bN|beEGi=(aN) Hier.)

\end{note}

\begin{note}[图片 OCR 摘录，待人工复核]
以下内容来自源图片识别，便于审核时对照：

\textbf{图片 1}（image.png）
\begin{Verbatim}[fontsize=\small]
5) OBEN E— FAR CB, SURG RE tHe ACRE.

(aN)(bN) =(ab)N =(ba)N =(bN )(aN)
6) PEAR BEHVLE— FAAS EE, LETRA AL BE.

Ck G=(a) X87 N<G
APPA ARRH, ATARI THAN TH, We

NIG.
VbeG,b=a bN=a'N=(aN)’
mh 8 8G/N={bN|beEGi=(aN) Hier.)
\end{Verbatim}

\end{note}

\begin{note}[来源路径]

近世代数/近世代数/无标题/交换群，循环群与商群 1cad2efa3f4280229ecffe1b0fe37bab.md

\end{note}

\section*{剩余类环的定义与符号 也就是商环}

\begin{note}[来源上级主题]

8剩余类环同态与理想 (8\%E5\%89\%A9\%E4\%BD\%99\%E7\%B1\%BB\%E7\%8E\%AF\%E5\%90\%8C\%E6\%80\%81\%E4\%B8\%8E\%E7\%90\%86\%E6\%83\%B3\%201edd2efa3f42801295a9f3bc6dd211c2.md)

\end{note}

\begin{definition}[剩余类环的定义与符号 也就是商环]

• **quotient ring**

定义2.11商环

命题2.9商环是环 当且仅当I是理想

\end{definition}

\begin{note}[关联图片]
以下图片已纳入本条整理，当前版本优先依据 Markdown 文字整理，图片细节可在审核时对照：

1. image.png\\

\end{note}

\begin{note}[来源路径]

近世代数/近世代数/无标题/剩余类环的定义与符号 也就是商环 1efd2efa3f428062bb60fef4a28718a3.md

\end{note}

\section*{命题1.30循环群的生成元的幂指数和循环群阶数互素}

\begin{note}[来源上级主题]

1.3有限群 (1\%203\%E6\%9C\%89\%E9\%99\%90\%E7\%BE\%A4\%201ead2efa3f4280c39a49d8b0af9aec03.md)

\end{note}

\begin{proposition}[命题1.30循环群的生成元的幂指数和循环群阶数互素]

所以循环群的生成元不止一个

这里的1，5和6互素

\end{proposition}

\begin{note}[图片 OCR 摘录，待人工复核]
以下内容来自源图片识别，便于审核时对照：

\textbf{图片 1}（IMG\_2976.jpeg）
\begin{Verbatim}[fontsize=\small]
fiz Si 1.30
& G = (x) Rn WMA. PGR Lm <n, Rx” A
my _   n
bl = ccd(nm)                (1.50)
a
WEY) BRL <m <n-1, Ril BRA EH REE oF ax Se. HF [xlan, HIE F nimk.
EPAARMNZAARE HOSA. CHAK nfm HALA, RNA
n    mig                 1.51
gcd(n, m)|gcd(n, m)                             Gen)
7 AA Salata 7a dG AKIN, UKE FSENT
n
=a                   (1.52)
LBL, ANS ER KER Sat RHA T ER.
Le AT ELBE HEHE HEA BEY AE BTC PC.
fiz 1.31
& G = (x) eA n UMA, Wx" Sm <n) ePAMA, HARYG
gcd(m,n) = 1                                (1.53)
TRAB bE b HAA, KALA MAA PALER O(n).                     “
\end{Verbatim}

\textbf{图片 2}（image.png）
\begin{Verbatim}[fontsize=\small]
- 0:

0,0,0,... > ARES 10}, ABATE

°1:

1,2,3,4,5,0 > S2SBrk > BEM

+2:                                             _
2,4,0,2,4,0,..> RA {0.2.4}, RBS > REAR
+3:

4:

° 5:

5,4,3,2,1,0 > S2/Se8 > BAM

ait:

Z, 15

. gcd(1,6) =1

. gcd(5,6) = 1

Prue (6) = 2 ME RR.
\end{Verbatim}

\end{note}

\begin{note}[来源路径]

近世代数/近世代数/无标题/命题1 30循环群的生成元的幂指数和循环群阶数互素 1d2d2efa3f4280cfaefbd295458b0a99.md

\end{note}

\section*{命题2.11交换环就能写成Ra}

\begin{proposition}[命题2.11交换环就能写成Ra]

fir jel 2.11

\& (R,+,) RV RAH, MaeR, Mi

(a) = Ra= \{ra:reéR\}                    (2.74)

—firh, Bai--:,aneR, Mi

(a1,°°+ ,dn) = Ra, +--+ +Ray = \{rja, +++ nan 2 71,°°* »Tn € R\}          (2.75) “

\end{proposition}

\begin{note}[图片 OCR 摘录，待人工复核]
以下内容来自源图片识别，便于审核时对照：

\textbf{图片 1}（image.png）
\begin{Verbatim}[fontsize=\small]
fir jel 2.11
& (R,+,) RV RAH, MaeR, Mi
(a) = Ra= {ra:reéR}                    (2.74)
—firh, Bai--:,aneR, Mi
(a1,°°+ ,dn) = Ra, +--+ +Ray = {rja, +++ nan 2 71,°°* »Tn € R}          (2.75) “
\end{Verbatim}

\end{note}

\begin{note}[来源路径]

近世代数/近世代数/无标题/命题2 11交换环就能写成Ra 209d2efa3f42800cab37e2d8d872bd5c.md

\end{note}

\section*{命题2.1负元 0元的运算}

\begin{note}[来源上级主题]

环 (\%E7\%8E\%AF\%201d9d2efa3f4280f38e94de87f834c519.md)

\end{note}

\begin{proposition}[命题2.1负元 0元的运算]

子级 项目: 证明 (\%E8\%AF\%81\%E6\%98\%8E\%201edd2efa3f4280e090e4d05710609507.md)

0是单位元

不写符号就是乘法

意思是：**任何元素乘以加法单位元 0 都等于 0**

乘法和 0 的关系建立在分配律基础上可证。

- b 是 b 的**加法逆元**，不是“减 b”

- 整个式子说明：**负数乘法的规律：乘法中取负，可以移到任何一个因子上**

- 在环中不需要交换律，这条等式依然成立（可用分配律证明）

\end{proposition}

\begin{note}[关联图片]
以下图片已纳入本条整理，当前版本优先依据 Markdown 文字整理，图片细节可在审核时对照：

1. image.png\\

\end{note}

\begin{note}[来源路径]

近世代数/近世代数/无标题/命题2 1负元 0元的运算 1d8d2efa3f4280408c76c66618c847b3.md

\end{note}

\section*{命题2.2 零环 加法单位元与乘法单位元相等}

\begin{note}[来源上级主题]

环 (\%E7\%8E\%AF\%201d9d2efa3f4280f38e94de87f834c519.md)

\end{note}

\begin{proposition}[命题2.2 零环 加法单位元与乘法单位元相等]

子级 项目: 证明 (\%E8\%AF\%81\%E6\%98\%8E\%201edd2efa3f428098a226ca9720be1aba.md)

\end{proposition}

\begin{note}[关联图片]
以下图片已纳入本条整理，当前版本优先依据 Markdown 文字整理，图片细节可在审核时对照：

1. 5507d41c-6a12-497f-a1d4-3ef5a52a3935.png\\
2. image.png\\

\end{note}

\begin{note}[来源路径]

近世代数/近世代数/无标题/命题2 2 零环 加法单位元与乘法单位元相等 1d8d2efa3f428072bc30db04dbb3356f.md

\end{note}

\section*{命题2.5环的直积还是环}

\begin{note}[来源上级主题]

环 (\%E7\%8E\%AF\%201d9d2efa3f4280f38e94de87f834c519.md)

\end{note}

\begin{proposition}[命题2.5环的直积还是环]

文本: 39

验证分配律

\end{proposition}

\begin{note}[图片 OCR 摘录，待人工复核]
以下内容来自源图片识别，便于审核时对照：

\textbf{图片 1}（image.png）
\begin{Verbatim}[fontsize=\small]
fir eh 2.5
iE)
TREE AE EAR IOM BR RAY, BW we [Te Ri MK RB, CREM RALH. Alb RA
BIER IE N AGO. RENE, RM REWADRE.
& (xdier Wier Zidier € Tie Ri, Wl
(xi)ier > (Nidier + (Za)ier)                                                (2.25)
H(i (Wi + Zi) ier                                                         (2.26)
=(x)ier > (Wiier + (Zaidier)                                                (2.27)
=(xiier (Vidier + (Miia * (Zidier                                        (2.28)
Ath, Tle Riet.:) 2-PbH. RRMIEW ThA.
\end{Verbatim}

\textbf{图片 2}（image 1.png）
\begin{Verbatim}[fontsize=\small]
» (2,3) + (5,7) = (7, 10)

« (2,3) - (5,7) = (10, 21)

RTH Z x ZHAMTSERH ZW, STR, BREE, ANCASAT, HIM:
(1,0) - (0, 1) = (0,0)
\end{Verbatim}

\end{note}

\begin{note}[来源路径]

近世代数/近世代数/无标题/命题2 5环的直积还是环 1d9d2efa3f42807da879eee76a3f0d7e.md

\end{note}

\section*{命题2.6交换环的理想左理想就是右理想}

\begin{note}[来源上级主题]

环 (\%E7\%8E\%AF\%201d9d2efa3f4280f38e94de87f834c519.md)

\end{note}

\begin{proposition}[命题2.6交换环的理想左理想就是右理想]

理想在乘法下吸收了整个环到理想上

\end{proposition}

\begin{note}[图片 OCR 摘录，待人工复核]
以下内容来自源图片识别，便于审核时对照：

\textbf{图片 1}（image.png）
\begin{Verbatim}[fontsize=\small]
fir jl 2.6
WRIK (Ri+,-) R-PREH, MI e-PARBMSARE I R-PSHM, LHARHIA-PERA.     “
\end{Verbatim}

\textbf{图片 2}（image 1.png）
\begin{Verbatim}[fontsize=\small]
WEN
PAE RAM RIE RRE, RERAM.
SAT VAL, FEISS AEE AE: FRREFIE PF “Wu” OE SEPSIS LL, te
RIcI                                                    (2.39)
IRCI                                                      (2.40)
HP FeWI, “4A ATCRIRRNEA
FAS ME BF nZ FEZ EE. HIRI, TRAE, ERE “WR” FESR
FHA SH.
\end{Verbatim}

\textbf{图片 3}（image 2.png）
\begin{Verbatim}[fontsize=\small]
ES GIF IR IRAA :

KR=Z, BAH

ype [= 2Z = {0,+2,+4,...}

+ ZAIBAR, BEBAR. EEA

- MHA, AB, RETBARTESH
Rime

(BH:

-2€1,3¢1, PA2-3=6¢1M

- (23: 3=9¢1X% = HAWES
& PACKET H
\end{Verbatim}

\textbf{图片 4}（image 3.png）
\begin{Verbatim}[fontsize=\small]
A BBS FHM RAE:

tt      FRRRE HARE
pssst   |     |
mE     |
WHA   |     | x
main Xam) xX

TBI (Ze) x    | (7278)
mB (SR) x    | GO (a7278)
\end{Verbatim}

\end{note}

\begin{note}[来源路径]

近世代数/近世代数/无标题/命题2 6交换环的理想左理想就是右理想 1edd2efa3f4280d6b552d742f9f0113f.md

\end{note}

\section*{命题2.7整环乘整数都是理想}

\begin{note}[来源上级主题]

环 (\%E7\%8E\%AF\%201d9d2efa3f4280f38e94de87f834c519.md)

\end{note}

\begin{proposition}[命题2.7整环乘整数都是理想]

因为一个是n的倍数的数再乘任意数都还是n的倍数

\end{proposition}

\begin{note}[图片 OCR 摘录，待人工复核]
以下内容来自源图片识别，便于审核时对照：

\textbf{图片 1}（image.png）
\begin{Verbatim}[fontsize=\small]
fir jel 2.7
nZ AZ                             (2.41)
4
WED) Be, Rae (nZ, +) &(Z4) BW (nk) FH. RAMA mh mn (WK) BARRA (K
Mk-#iEA).

HK, ERSZE-PREM, REN RAE RI Cl, EKEMEZ nZcnZ. TARR? RE, —
SERRE n WER, Renheh. REGAN, PRA LARRS HOY. RERMBRASE—
HRY, RATHRAEAEX. RERAN, FERAW AEH SLRMRRBRIR. KAT WAH RA RK
FWA, WRK BOWE.

BiEWZ-nZcnZ, RRMA me Z, nk enZ(keZ), REE mnk enZ WW. HKEAH

mnk =n(mk) €nZ                                  (2.42)

REAR, RMIT ZZ HBR.

FATAL, TERNIE REPEAT DAMPER AS He CRIA AS EK) , BEAT DAA TAS ee SC NASH) «
FRA] FANS | FHS I I PS SES
\end{Verbatim}

\end{note}

\begin{note}[来源路径]

近世代数/近世代数/无标题/命题2 7整环乘整数都是理想 1edd2efa3f4280879de0f548a3792784.md

\end{note}

\section*{命题2.8 R-R‘环同态 核都是理想 像都是子环}

\begin{note}[来源上级主题]

环 (\%E7\%8E\%AF\%201d9d2efa3f4280f38e94de87f834c519.md)

\end{note}

\begin{proposition}[命题2.8 R-R‘环同态 核都是理想 像都是子环]

子级 项目: 证明 (\%E8\%AF\%81\%E6\%98\%8E\%201edd2efa3f428054a04be17abb7ec9aa.md)

定义2.11商环

\end{proposition}

\begin{note}[关联图片]
以下图片已纳入本条整理，当前版本优先依据 Markdown 文字整理，图片细节可在审核时对照：

1. image.png\\

\end{note}

\begin{note}[来源路径]

近世代数/近世代数/无标题/命题2 8 R-R‘环同态 核都是理想 像都是子环 1edd2efa3f4280c7aab3cfeaab813a46.md

\end{note}

\section*{命题2.9商环是环 当且仅当I是理想}

\begin{note}[来源上级主题]

环 (\%E7\%8E\%AF\%201d9d2efa3f4280f38e94de87f834c519.md)

\end{note}

\begin{proposition}[命题2.9商环是环 当且仅当I是理想]

子级 项目: 证明 (\%E8\%AF\%81\%E6\%98\%8E\%201eed2efa3f42807fbefeddb9976e5a4c.md)

文本: 42

\end{proposition}

\begin{note}[关联图片]
以下图片已纳入本条整理，当前版本优先依据 Markdown 文字整理，图片细节可在审核时对照：

1. image.png\\

\end{note}

\begin{note}[来源路径]

近世代数/近世代数/无标题/命题2 9商环是环 当且仅当I是理想 1eed2efa3f4280c88023ce1caa85f335.md

\end{note}

\section*{域就是交换的除环}

\begin{note}[来源上级主题]

第三章 (\%E7\%AC\%AC\%E4\%B8\%89\%E7\%AB\%A0\%201ecd2efa3f4280c7a6b0ea0de416c09f.md)

\end{note}

\begin{note}[域就是交换的除环]

定义2.5域交换的除环——更好的环

\end{note}

\begin{note}[来源路径]

近世代数/近世代数/无标题/域就是交换的除环 1edd2efa3f428082b256f54c7b69c4d8.md

\end{note}

\section*{奇偶排列素数阶的群一定是循环群}

\begin{note}[来源上级主题]

子群的陪集 (\%E5\%AD\%90\%E7\%BE\%A4\%E7\%9A\%84\%E9\%99\%AA\%E9\%9B\%86\%201c3d2efa3f42806a85b0f0113028b207.md)

\end{note}

\begin{note}[奇偶排列素数阶的群一定是循环群]

待根据图片进一步人工校对。

\end{note}

\begin{note}[来源路径]

近世代数/近世代数/无标题/奇偶排列素数阶的群一定是循环群 1bdd2efa3f428081b6cdc32b80a5f208.md

\end{note}

\section*{子环与子除环}

\begin{note}[来源上级主题]

第三章 (\%E7\%AC\%AC\%E4\%B8\%89\%E7\%AB\%A0\%201ecd2efa3f4280c7a6b0ea0de416c09f.md)

\end{note}

\begin{note}[子环与子除环]

EM —PHRVN-TRFES WMRN-TbFR, BHSABUIPR WV

RBGE ARLE R— TH.

—SRARN—-TRFRS WRR N—-TSFRR, BH SABI R OK

BUS EK VE KTR.

Ay, RMT UM EF BR, FN Ras.

— HBAS FR S TEMPS FH BA RE

a, bES=>a—-bES, abeS.

—TS RAB —RFR S TER—DF RA FEA AE

QD SAB—-tTAS FSI;

Gi)               a, bE S>a—-bes,

a, bES, bA0=>ab"ES,

\end{note}

\begin{note}[图片 OCR 摘录，待人工复核]
以下内容来自源图片识别，便于审核时对照：

\textbf{图片 1}（image.png）
\begin{Verbatim}[fontsize=\small]
EM —PHRVN-TRFES WMRN-TbFR, BHSABUIPR WV
RBGE ARLE R— TH.
—SRARN—-TRFRS WRR N—-TSFRR, BH SABI R OK
BUS EK VE KTR.
Ay, RMT UM EF BR, FN Ras.
— HBAS FR S TEMPS FH BA RE
a, bES=>a—-bES, abeS.
—TS RAB —RFR S TER—DF RA FEA AE
QD SAB—-tTAS FSI;
Gi)               a, bE S>a—-bes,
a, bES, bA0=>ab"ES,
\end{Verbatim}

\end{note}

\begin{note}[来源路径]

近世代数/近世代数/无标题/子环与子除环 1edd2efa3f42803d86f8f66dab16dc23.md

\end{note}

\section*{子环的等价条件引理2.1}

\begin{note}[来源上级主题]

定义2.6子环 群是子群 幺半群是子幺半群 符号< (\%E5\%AE\%9A\%E4\%B9\%892\%206\%E5\%AD\%90\%E7\%8E\%AF\%20\%E7\%BE\%A4\%E6\%98\%AF\%E5\%AD\%90\%E7\%BE\%A4\%20\%E5\%B9\%BA\%E5\%8D\%8A\%E7\%BE\%A4\%E6\%98\%AF\%E5\%AD\%90\%E5\%B9\%BA\%E5\%8D\%8A\%E7\%BE\%A4\%20\%E7\%AC\%A6\%E5\%8F\%B7\%201d8d2efa3f42800a807fc311bee30d05.md)

\end{note}

\begin{lemma}[子环的等价条件引理2.1]

a-b蕴含着逆元

\end{lemma}

\begin{note}[图片 OCR 摘录，待人工复核]
以下内容来自源图片识别，便于审核时对照：

\textbf{图片 1}（image.png）
\begin{Verbatim}[fontsize=\small]
Sec, DemtAL UE (S,+) 2 (CR, +) PRE, (S,-) 2 (R) PAAR. BRAT Ta RE
5 |B 2.1
A (Ri4+,) EPH, MSCR, MS<RHARS
les                             (2.20)
Va,béeS,a-—b,abeS                                       (2.21)
9
UE Rae RT ASAE, MAO=1-1€S. Hh -a=0-aeS, a+b=a-(-b) ES. RHREA TRE
FR.
A-VAAELAW. AMS ETH, MAa-b=a+(-b)eS.
PANO Z iE Q HF.
AS FH, FATALE AE ATT DEM TE. ARIE REN. ANEMIA A EPPA Fab
WASERTE TT, DAB ANG T
\end{Verbatim}

\end{note}

\begin{note}[来源路径]

近世代数/近世代数/无标题/子环的等价条件引理2 1 1edd2efa3f4280d79067f6538b96b7ec.md

\end{note}

\section*{定义1.25循环群}

\begin{note}[来源上级主题]

1.3有限群 (1\%203\%E6\%9C\%89\%E9\%99\%90\%E7\%BE\%A4\%201ead2efa3f4280c39a49d8b0af9aec03.md)

\end{note}

\begin{definition}[定义1.25循环群]

\& (G+) RA GH EXE G, (RAG G = (x), RG RARA—- MAR, ta x HARA GAY ERA. *

\end{definition}

\begin{note}[图片 OCR 摘录，待人工复核]
以下内容来自源图片识别，便于审核时对照：

\textbf{图片 1}（IMG\_2950.jpeg）
\begin{Verbatim}[fontsize=\small]
& (G+) RA GH EXE G, (RAG G = (x), RG RARA—- MAR, ta x HARA GAY ERA. *
\end{Verbatim}

\textbf{图片 2}（IMG\_2951.jpeg）
\begin{Verbatim}[fontsize=\small]
FEx 2
2 7 ——  SS
wa                  x
x4
1                        e
.                  ys
‘SN             x
HE x      FE x      FE x      He x      HE x
Fe ce Le 6 me pr’ ------p
Pe] 1.3: ABR AU FC BR A HE.
FREEWAY, th
\end{Verbatim}

\end{note}

\begin{note}[来源路径]

近世代数/近世代数/无标题/定义1 25循环群 1d1d2efa3f42803f81baf1a7f0cbd190.md

\end{note}

\section*{定义2.10理想}

\begin{definition}[定义2.10理想]

EL 2.10

S(R+:) Z-PH, MICR, MIL, HIRRYUABM, SF

(I, +) < (R, +)                                                (2.34)

Vr é€R,Va€L,ra€1                                          (2.35)

Ree, SUTRA RYH, S

(I, +) < (R, +)                                                (2.36)

Va €1,Vr € R,ar €1                                          (2.37)

toR PR APRMLRABM, AHA RR APL, ET AR.                    .

SIMA Heat, SEE IRA ETH. Sb, —-MHRUE TMA M4E BEM.

\end{definition}

\begin{note}[图片 OCR 摘录，待人工复核]
以下内容来自源图片识别，便于审核时对照：

\textbf{图片 1}（image.png）
\begin{Verbatim}[fontsize=\small]
EL 2.10
S(R+:) Z-PH, MICR, MIL, HIRRYUABM, SF
(I, +) < (R, +)                                                (2.34)
Vr é€R,Va€L,ra€1                                          (2.35)
Ree, SUTRA RYH, S
(I, +) < (R, +)                                                (2.36)
Va €1,Vr € R,ar €1                                          (2.37)
toR PR APRMLRABM, AHA RR APL, ET AR.                    .
SIMA Heat, SEE IRA ETH. Sb, —-MHRUE TMA M4E BEM.
\end{Verbatim}

\end{note}

\begin{note}[来源路径]

近世代数/近世代数/无标题/定义2 10理想 1d9d2efa3f428084bdbade477843818e.md

\end{note}

\section*{定义2.10理想 理想的符号也是▶️指向理想}

\begin{note}[来源上级主题]

环 (\%E7\%8E\%AF\%201d9d2efa3f4280f38e94de87f834c519.md)

\end{note}

\begin{definition}[定义2.10理想 理想的符号也是▶️指向理想]

子级 项目: 引理2.2理想是子环当且仅当理想是正规环 (\%E5\%BC\%95\%E7\%90\%862\%202\%E7\%90\%86\%E6\%83\%B3\%E6\%98\%AF\%E5\%AD\%90\%E7\%8E\%AF\%E5\%BD\%93\%E4\%B8\%94\%E4\%BB\%85\%E5\%BD\%93\%E7\%90\%86\%E6\%83\%B3\%E6\%98\%AF\%E6\%AD\%A3\%E8\%A7\%84\%E7\%8E\%AF\%201edd2efa3f4280f38dbaeb87e2cd26c6.md), 无标题 (\%E6\%97\%A0\%E6\%A0\%87\%E9\%A2\%98\%201edd2efa3f4280e7b9eedcae01414587.md)

(I,+)是子群，左乘环中任意元素还在环里面就是左理想同理右理想

底下那句话不对，理想都是子环

\end{definition}

\begin{note}[关联图片]
以下图片已纳入本条整理，当前版本优先依据 Markdown 文字整理，图片细节可在审核时对照：

1. image.png\\

\end{note}

\begin{note}[来源路径]

近世代数/近世代数/无标题/定义2 10理想 理想的符号也是▶️指向理想 1edd2efa3f4280cbb2a2d6978eb4e340.md

\end{note}

\section*{定义2.11商环}

\begin{note}[来源上级主题]

环 (\%E7\%8E\%AF\%201d9d2efa3f4280f38e94de87f834c519.md)

\end{note}

\begin{definition}[定义2.11商环]

R是一个环 I是R的一个理想 商环就是R中元素对I做加法形成的陪集

a+I是一个集合

商环，商群是加法的配集，商环还是加法的陪集

\end{definition}

\begin{note}[关联图片]
以下图片已纳入本条整理，当前版本优先依据 Markdown 文字整理，图片细节可在审核时对照：

1. image.png\\
2. image 1.png\\

\end{note}

\begin{note}[来源路径]

近世代数/近世代数/无标题/定义2 11商环 1eed2efa3f4280d08291e336ed049bc4.md

\end{note}

\section*{定义2.13生成理想的定义}

\begin{definition}[定义2.13生成理想的定义]

所有包含他的理想的交集

定义1.2.4子集生成群证明

生成的理想是包含(4)的理想的交集

\end{definition}

\begin{note}[关联图片]
以下图片已纳入本条整理，当前版本优先依据 Markdown 文字整理，图片细节可在审核时对照：

1. image.png\\
2. image 1.png\\

\end{note}

\begin{note}[来源路径]

近世代数/近世代数/无标题/定义2 13生成理想的定义 209d2efa3f4280cfbf73e662a3d86e2f.md

\end{note}

\section*{定义2.14 主理想}

\begin{definition}[定义2.14 主理想]

EM 2.14

(a) = (\{a\})                                                (2.72)

Aw a EA EPL. AR, BPM AERP AAEM, AMBALA,

MT aan ER, RAVEN

(a1,°-- ,@n) = (\{a1,-- ,an\})                                  (2.73)

Ait, BrP PLM ARAIRA THERM, RANA A AIRE MY.

\&

\end{definition}

\begin{note}[图片 OCR 摘录，待人工复核]
以下内容来自源图片识别，便于审核时对照：

\textbf{图片 1}（image.png）
\begin{Verbatim}[fontsize=\small]
EM 2.14
(a) = ({a})                                                (2.72)
Aw a EA EPL. AR, BPM AERP AAEM, AMBALA,
MT aan ER, RAVEN
(a1,°-- ,@n) = ({a1,-- ,an})                                  (2.73)
Ait, BrP PLM ARAIRA THERM, RANA A AIRE MY.
&
\end{Verbatim}

\end{note}

\begin{note}[来源路径]

近世代数/近世代数/无标题/定义2 14 主理想 209d2efa3f428074a610e5ed31f4a5cb.md

\end{note}

\section*{定义2.2交换环}

\begin{note}[来源上级主题]

环 (\%E7\%8E\%AF\%201d9d2efa3f4280f38e94de87f834c519.md)

\end{note}

\begin{definition}[定义2.2交换环]

交换环是给乘法补上了交换率，n*n实矩阵是一个非交换环

\end{definition}

\begin{note}[关联图片]
以下图片已纳入本条整理，当前版本优先依据 Markdown 文字整理，图片细节可在审核时对照：

1. image.png\\

\end{note}

\begin{note}[来源路径]

近世代数/近世代数/无标题/定义2 2交换环 1d8d2efa3f4280d4b681c6f9125a7bd9.md

\end{note}

\section*{定义2.3 有乘法可逆元素叫做单位除环的性质很好}

\begin{note}[来源上级主题]

环 (\%E7\%8E\%AF\%201d9d2efa3f4280f38e94de87f834c519.md)

\end{note}

\begin{definition}[定义2.3 有乘法可逆元素叫做单位除环的性质很好]

引理1.1幺半群的逆元构成群

\end{definition}

\begin{note}[图片 OCR 摘录，待人工复核]
以下内容来自源图片识别，便于审核时对照：

\textbf{图片 1}（image.png）
\begin{Verbatim}[fontsize=\small]
ESL 2.3
? Ay B45.
*
FALE, ARN (ARASH), OF 1, ALO AA FAWOY , ROA a € R,
FIVE Oa = O04 1. ARDS AITR OL AT Regie R\ {0} = RX, BUTAIRS TCR ADA AGA WOT. LORRI ERA EE te
REN, RMKEA-TRF.
\end{Verbatim}

\end{note}

\begin{note}[来源路径]

近世代数/近世代数/无标题/定义2 3 有乘法可逆元素叫做单位除环的性质很好 1d8d2efa3f4280ae9d4bc1dfa10ee023.md

\end{note}

\section*{定义2.4除环非零元素都是单位}

\begin{note}[来源上级主题]

环 (\%E7\%8E\%AF\%201d9d2efa3f4280f38e94de87f834c519.md)

\end{note}

\begin{definition}[定义2.4除环非零元素都是单位]

除环非零元素都有逆元

加法是阿贝尔群，乘法构成群而不是幺半群 乘法对加法有左右分配率

证明一个环是除环也就是任意两个数乘起来只要两个都不是0结果就不能是0

证明域就是证明 除环的任意元素乘法都可以交换

\end{definition}

\begin{note}[关联图片]
以下图片已纳入本条整理，当前版本优先依据 Markdown 文字整理，图片细节可在审核时对照：

1. image.png\\

\end{note}

\begin{note}[来源路径]

近世代数/近世代数/无标题/定义2 4除环非零元素都是单位 1d8d2efa3f428059bd92f93a477bc7b1.md

\end{note}

\section*{定义2.5域交换的除环——更好的环}

\begin{note}[来源上级主题]

环 (\%E7\%8E\%AF\%201d9d2efa3f4280f38e94de87f834c519.md)

\end{note}

\begin{definition}[定义2.5域交换的除环——更好的环]

乘法对加法有分配律别忘了

\end{definition}

\begin{note}[图片 OCR 摘录，待人工复核]
以下内容来自源图片识别，便于审核时对照：

\textbf{图片 1}（image.png）
\begin{Verbatim}[fontsize=\small]
fi jh 2.3
(R,+,°) e—-MK, SARS
(R,+) AM WA RAF                                    (2.14)
(R \ {0}, +) RAPT RAE                                   (2.15)
KLM AKA DAE                                    (2.16)
a
\end{Verbatim}

\end{note}

\begin{note}[来源路径]

近世代数/近世代数/无标题/定义2 5域交换的除环——更好的环 1d8d2efa3f4280daabf6fdeddf87f1ff.md

\end{note}

\section*{定义2.6子环 群是子群 幺半群是子幺半群 符号<}

\begin{note}[来源上级主题]

环 (\%E7\%8E\%AF\%201d9d2efa3f4280f38e94de87f834c519.md)

\end{note}

\begin{definition}[定义2.6子环 群是子群 幺半群是子幺半群 符号<]

子级 项目: 子环的等价条件引理2.1 (\%E5\%AD\%90\%E7\%8E\%AF\%E7\%9A\%84\%E7\%AD\%89\%E4\%BB\%B7\%E6\%9D\%A1\%E4\%BB\%B6\%E5\%BC\%95\%E7\%90\%862\%201\%201edd2efa3f4280d79067f6538b96b7ec.md)

有加法单位元，乘法单位元 封闭 加群是父群的子群，乘幺半群是父幺半群的子幺半群

\end{definition}

\begin{note}[关联图片]
以下图片已纳入本条整理，当前版本优先依据 Markdown 文字整理，图片细节可在审核时对照：

1. image.png\\

\end{note}

\begin{note}[来源路径]

近世代数/近世代数/无标题/定义2 6子环 群是子群 幺半群是子幺半群 符号 1d8d2efa3f42800a807fc311bee30d05.md

\end{note}

\section*{定义2.7子集A生成环<A>}

\begin{note}[来源上级主题]

环 (\%E7\%8E\%AF\%201d9d2efa3f4280f38e94de87f834c519.md)

\end{note}

\begin{definition}[定义2.7子集A生成环<A>]

\& (R,4+,:) E-PH, MACR, MABRMATH, ite (A), LLAMAAST AMTFHHRAK, BP

(A) =( \{SC R:S3 4,5 <R\}                    (2.22)

\&

\end{definition}

\begin{note}[图片 OCR 摘录，待人工复核]
以下内容来自源图片识别，便于审核时对照：

\textbf{图片 1}（image.png）
\begin{Verbatim}[fontsize=\small]
& (R,4+,:) E-PH, MACR, MABRMATH, ite (A), LLAMAAST AMTFHHRAK, BP
(A) =( {SC R:S3 4,5 <R}                    (2.22)
&
\end{Verbatim}

\textbf{图片 2}（image 1.png）
\begin{Verbatim}[fontsize=\small]
ari 2.4                           —
VED) BARSRRAEEH, AARAAHRE-MAET ANTH.
HPA, RNAALH HTH. SSHe-PASTANTR. AALEB-PREN SH, P(A).
ane A), illa—b,ab £e—-SREH SH, AAB—-DSS BETH. Alb a—b,ab (A).
oe Le, (A) < R.
Vise FFM, FTAA. AACE. TRE, PREG Ma BEEBE)
\end{Verbatim}

\end{note}

\begin{note}[来源路径]

近世代数/近世代数/无标题/定义2 7子集A生成环 A 1d8d2efa3f4280f9b734fb70e49c8370.md

\end{note}

\section*{定义2.8环的直积}

\begin{note}[来源上级主题]

环 (\%E7\%8E\%AF\%201d9d2efa3f4280f38e94de87f834c519.md)

\end{note}

\begin{definition}[定义2.8环的直积]

一族环 多个环

\end{definition}

\begin{note}[图片 OCR 摘录，待人工复核]
以下内容来自源图片识别，便于审核时对照：

\textbf{图片 1}（image.png）
\begin{Verbatim}[fontsize=\small]
& (Ri, tis diel RHA. AME LEMH HR, A (jer Ri. t.-)o SF Wadier, ier € [er Ri, HA
\end{Verbatim}

\textbf{图片 2}（image 1.png）
\begin{Verbatim}[fontsize=\small]
eo
(xiier + Oiier = i +i Vidier                                       (2.23)
(xiier - (Vidier = (Xi it Yadier                                     (2.24)
&
\end{Verbatim}

\end{note}

\begin{note}[来源路径]

近世代数/近世代数/无标题/定义2 8环的直积 1d9d2efa3f42800ea898d0b403b41d0d.md

\end{note}

\section*{定义2.9环同态}

\begin{note}[来源上级主题]

环 (\%E7\%8E\%AF\%201d9d2efa3f4280f38e94de87f834c519.md)

\end{note}

\begin{definition}[定义2.9环同态]

文本: 39

定义1.15群同态下定义 核是映射到e‘的原像集 imf是像集

定义1.5幺半群的同态 要保持单位元

乘法单位元传递

加法乘法能拆括号

\end{definition}

\begin{note}[关联图片]
以下图片已纳入本条整理，当前版本优先依据 Markdown 文字整理，图片细节可在审核时对照：

1. image.png\\
2. image 1.png\\

\end{note}

\begin{note}[来源路径]

近世代数/近世代数/无标题/定义2 9环同态 1d9d2efa3f428053a89bd4f291079c1e.md

\end{note}

\section*{定理1 理想和正规子群很像 模理想的剩余类还是环而且和原环同态}

\begin{note}[来源上级主题]

8剩余类环同态与理想 (8\%E5\%89\%A9\%E4\%BD\%99\%E7\%B1\%BB\%E7\%8E\%AF\%E5\%90\%8C\%E6\%80\%81\%E4\%B8\%8E\%E7\%90\%86\%E6\%83\%B3\%201edd2efa3f42801295a9f3bc6dd211c2.md)

\end{note}

\begin{theorem}[定理1 理想和正规子群很像 模理想的剩余类还是环而且和原环同态]

EH1 PERE-TH, VREWN-THA, RAMARUHMRKE

RHR. BMARARHE-TH, FAR GR AA.

° 88 +

\end{theorem}

\begin{note}[图片 OCR 摘录，待人工复核]
以下内容来自源图片识别，便于审核时对照：

\textbf{图片 1}（image.png）
\begin{Verbatim}[fontsize=\small]
EH1 PERE-TH, VREWN-THA, RAMARUHMRKE
RHR. BMARARHE-TH, FAR GR AA.
° 88 +
\end{Verbatim}

\textbf{图片 2}（image 1.png）
\begin{Verbatim}[fontsize=\small]
WEAR BR
a— [a]
GREER BIR H-*+HABH,. PURER AA, WR E—-TPH. iE.
EM RUBRAR HRAUHMRER. RPRHRNAAS
R/X
RRM.
\end{Verbatim}

\end{note}

\begin{note}[来源路径]

近世代数/近世代数/无标题/定理1 理想和正规子群很像 模理想的剩余类还是环而且和原环同态 1edd2efa3f428071ac13fae2aef057df.md

\end{note}

\section*{定理1除环只有两个理想0理想和单位理想}

\begin{note}[来源上级主题]

7理想 (7\%E7\%90\%86\%E6\%83\%B3\%201edd2efa3f4280a88bfce32f4c292dfb.md)

\end{note}

\begin{theorem}[定理1除环只有两个理想0理想和单位理想]

子级 项目: 无标题 (\%E6\%97\%A0\%E6\%A0\%87\%E9\%A2\%98\%201edd2efa3f42809db73ce9e5011f76a6.md)

交换的除环是域，所以域也只有两个理想零理想和单位理想

\end{theorem}

\begin{note}[关联图片]
以下图片已纳入本条整理，当前版本优先依据 Markdown 文字整理，图片细节可在审核时对照：

1. image.png\\

\end{note}

\begin{note}[来源路径]

近世代数/近世代数/无标题/定理1除环只有两个理想0理想和单位理想 1edd2efa3f4280198c05df2fd987f2b8.md

\end{note}

\section*{定理2 R和自己的商环同态 同态满射核就是理想}

\begin{note}[来源上级主题]

8剩余类环同态与理想 (8\%E5\%89\%A9\%E4\%BD\%99\%E7\%B1\%BB\%E7\%8E\%AF\%E5\%90\%8C\%E6\%80\%81\%E4\%B8\%8E\%E7\%90\%86\%E6\%83\%B3\%201edd2efa3f42801295a9f3bc6dd211c2.md)

\end{note}

\begin{theorem}[定理2 R和自己的商环同态 同态满射核就是理想]

命题2.8 R-R‘环同态 核都是理想 像都是子环

命题2.9商环是环 当且仅当I是理想

定义1.15群同态下定义 核是映射到e‘的原像集 imf是像集

核在环里面是映射到0元

\end{theorem}

\begin{note}[关联图片]
以下图片已纳入本条整理，当前版本优先依据 Markdown 文字整理，图片细节可在审核时对照：

1. image.png\\

\end{note}

\begin{note}[来源路径]

近世代数/近世代数/无标题/定理2 R和自己的商环同态 同态满射核就是理想 1edd2efa3f4280868a83cd8270b657a6.md

\end{note}

\section*{定理2 n元置换可以写成循环置换的乘积}

\begin{note}[来源上级主题]

循环置换  (\%E5\%BE\%AA\%E7\%8E\%AF\%E7\%BD\%AE\%E6\%8D\%A2\%201bad2efa3f4280109703cec9c3fb3605.md)

\end{note}

\begin{theorem}[定理2 n元置换可以写成循环置换的乘积]

22 B—-Phn SUN RR MTU SMATTEARAR EAS

CA FEE HY) VE BRAY FER.

\end{theorem}

\begin{note}[图片 OCR 摘录，待人工复核]
以下内容来自源图片识别，便于审核时对照：

\textbf{图片 1}（image.png）
\begin{Verbatim}[fontsize=\small]
22 B—-Phn SUN RR MTU SMATTEARAR EAS
CA FEE HY) VE BRAY FER.
\end{Verbatim}

\end{note}

\begin{note}[来源路径]

近世代数/近世代数/无标题/定理2 n元置换可以写成循环置换的乘积 1bad2efa3f428028af82e07c1d9e66d5.md

\end{note}

\section*{定理2 整环，除环，域特征只能是素数或者无限}

\begin{note}[来源上级主题]

第三章 (\%E7\%AC\%AC\%E4\%B8\%89\%E7\%AB\%A0\%201ecd2efa3f4280c7a6b0ea0de416c09f.md)

\end{note}

\begin{theorem}[定理2 整环，除环，域特征只能是素数或者无限]

整环，除环，域都是没有零因子的

1.10阿贝尔群

\end{theorem}

\begin{note}[图片 OCR 摘录，待人工复核]
以下内容来自源图片识别，便于审核时对照：

\textbf{图片 1}（image.png）
\begin{Verbatim}[fontsize=\small]
EH 2 MRASATH RR MREBARBMAn, BAn k-TRK.
SBA IC a HU »
n,aX0, neaX0, {A (nia) (na) = (nj n2)a*=0,
x5R RASA FT ABE HS. UESE.
Hit EM. RAURRH ERE, ME—-TPRE p.
FE—PERMEE »p WARM E, A-ARABWHREK, ME
(a+b)? =a’+b?.
XxeAW
PY at        2    p-1_1 pp
(a+b)? =a? +(" a?  bte+( "dab  +0,
Ti(”) Bp By— AMSA SR.
\end{Verbatim}

\end{note}

\begin{note}[来源路径]

近世代数/近世代数/无标题/定理2 整环，除环，域特征只能是素数或者无限 1edd2efa3f42802dbf6fd904ddc8051a.md

\end{note}

\section*{定理2环同态传递零元负元 如果是交换环传递单位元}

\begin{note}[来源上级主题]

第三章 (\%E7\%AC\%AC\%E4\%B8\%89\%E7\%AB\%A0\%201ecd2efa3f4280c7a6b0ea0de416c09f.md)

\end{note}

\begin{theorem}[定理2环同态传递零元负元 如果是交换环传递单位元]

EH2 RERAR BATH. HFHRSERAA. WA. RHSrHR

BR WMS5c, R Moca hhh eHo Hes. HA. fee R EEK »

PAR RARARYH;, RRA 1. PAR HAMM 1. WAIHI

KR.

RINBER., —bHARARSAFIX-—+ERA ST —T+HMAW NEA

SHWRS. RMB.

\end{theorem}

\begin{note}[图片 OCR 摘录，待人工复核]
以下内容来自源图片识别，便于审核时对照：

\textbf{图片 1}（image.png）
\begin{Verbatim}[fontsize=\small]
EH2 RERAR BATH. HFHRSERAA. WA. RHSrHR
BR WMS5c, R Moca hhh eHo Hes. HA. fee R EEK »
PAR RARARYH;, RRA 1. PAR HAMM 1. WAIHI
KR.

RINBER., —bHARARSAFIX-—+ERA ST —T+HMAW NEA
SHWRS. RMB.
\end{Verbatim}

\end{note}

\begin{note}[来源路径]

近世代数/近世代数/无标题/定理2环同态传递零元负元 如果是交换环传递单位元 1edd2efa3f42805cb845f45123227986.md

\end{note}

\section*{定理3 R和R的商环同态满射的性质 传递}

\begin{note}[来源上级主题]

8剩余类环同态与理想 (8\%E5\%89\%A9\%E4\%BD\%99\%E7\%B1\%BB\%E7\%8E\%AF\%E5\%90\%8C\%E6\%80\%81\%E4\%B8\%8E\%E7\%90\%86\%E6\%83\%B3\%201edd2efa3f42801295a9f3bc6dd211c2.md)

\end{note}

\begin{theorem}[定理3 R和R的商环同态满射的性质 传递]

文本: 85

同态满射

\end{theorem}

\begin{note}[图片 OCR 摘录，待人工复核]
以下内容来自源图片识别，便于审核时对照：

\textbf{图片 1}（image.png）
\begin{Verbatim}[fontsize=\small]
EH3 FKR AAR H-tT+AMABHSE,
GQ) RW—-+*+FHS HRS BR BW—tT+FH;
Gi) R W—7A Be UW BILE R WT
Gii) R N—-*+FHS HRRS BRY-T+FH;
(iv) R AO? UY ROU ER WS.
\end{Verbatim}

\end{note}

\begin{note}[来源路径]

近世代数/近世代数/无标题/定理3 R和R的商环同态满射的性质 传递 1edd2efa3f4280a6a5a6cd9145845beb.md

\end{note}

\section*{定理3同构传递整环除环域的性质}

\begin{note}[来源上级主题]

第三章 (\%E7\%AC\%AC\%E4\%B8\%89\%E7\%AB\%A0\%201ecd2efa3f4280c7a6b0ea0de416c09f.md)

\end{note}

\begin{theorem}[定理3同构传递整环除环域的性质]

EH3 BRERMR BASH, FHARER. RA, GRABH, RB

SH; R ERM, RHERM; RE, REM.

FUP HWE, RINAN GRE -TH, ERERA—-THENH. Bie

BUX MIE Ne, BRAT Ae EF BP A RF eB 4 FRAT SE BH

5) heEERA ASA ZHFE—-T——EH go. HAA AMMAR

AK. MARNIE A Me MeMRK, HEA GA MEM RM

Fe YE Te i HE fal FY.

\end{theorem}

\begin{note}[图片 OCR 摘录，待人工复核]
以下内容来自源图片识别，便于审核时对照：

\textbf{图片 1}（image.png）
\begin{Verbatim}[fontsize=\small]
EH3 BRERMR BASH, FHARER. RA, GRABH, RB
SH; R ERM, RHERM; RE, REM.

FUP HWE, RINAN GRE -TH, ERERA—-THENH. Bie
BUX MIE Ne, BRAT Ae EF BP A RF eB 4 FRAT SE BH

5) heEERA ASA ZHFE—-T——EH go. HAA AMMAR
AK. MARNIE A Me MeMRK, HEA GA MEM RM
Fe YE Te i HE fal FY.
\end{Verbatim}

\end{note}

\begin{note}[来源路径]

近世代数/近世代数/无标题/定理3同构传递整环除环域的性质 1edd2efa3f4280ea855efe8275531d49.md

\end{note}

\section*{引理2 交换环只有零理想和单位理想R是一个域}

\begin{note}[来源上级主题]

9极大理想 (9\%E6\%9E\%81\%E5\%A4\%A7\%E7\%90\%86\%E6\%83\%B3\%201efd2efa3f42802fa98fcbcc6a6ed8b2.md)

\end{note}

\begin{lemma}[引理2 交换环只有零理想和单位理想R是一个域]

子级 项目: 证明 (\%E8\%AF\%81\%E6\%98\%8E\%20202d2efa3f4280ba8435d49791bf53ee.md)

交换环+只有零理想和单位理想就是一个域

\end{lemma}

\begin{note}[关联图片]
以下图片已纳入本条整理，当前版本优先依据 Markdown 文字整理，图片细节可在审核时对照：

1. image.png\\

\end{note}

\begin{note}[来源路径]

近世代数/近世代数/无标题/引理2 交换环只有零理想和单位理想R是一个域 202d2efa3f42801f8790d49e461b4508.md

\end{note}

\section*{引理2.2理想是子环当且仅当理想是正规环}

\begin{note}[来源上级主题]

定义2.10理想 理想的符号也是▶️指向理想 (\%E5\%AE\%9A\%E4\%B9\%892\%2010\%E7\%90\%86\%E6\%83\%B3\%20\%E7\%90\%86\%E6\%83\%B3\%E7\%9A\%84\%E7\%AC\%A6\%E5\%8F\%B7\%E4\%B9\%9F\%E6\%98\%AF\%E2\%96\%B6\%EF\%B8\%8F\%E6\%8C\%87\%E5\%90\%91\%E7\%90\%86\%E6\%83\%B3\%201edd2efa3f4280cbb2a2d6978eb4e340.md)

\end{note}

\begin{lemma}[引理2.2理想是子环当且仅当理想是正规环]

I是R的理想 I是R的子环当且仅当I=R 1是I里面的I里元素左乘R里面任何元素都在I里面

I是

\end{lemma}

\begin{note}[关联图片]
以下图片已纳入本条整理，当前版本优先依据 Markdown 文字整理，图片细节可在审核时对照：

1. image.png\\

\end{note}

\begin{note}[来源路径]

近世代数/近世代数/无标题/引理2 2理想是子环当且仅当理想是正规环 1edd2efa3f4280f38dbaeb87e2cd26c6.md

\end{note}

\section*{引理2.3从理想可以从环同态中定义}

\begin{note}[来源上级主题]

环 (\%E7\%8E\%AF\%201d9d2efa3f4280f38e94de87f834c519.md)

\end{note}

\begin{lemma}[引理2.3从理想可以从环同态中定义]

文本: 41

定义1.15群同态下定义 核是映射到e‘的原像集 imf是像集

ker(f)=e’也就是nZ 比如说f(1)=1+nZ

Zn是模n的剩余类加群

因为f(ker(f))=0 所以ker(f)=nZ是R的理想

加法的群同态已经证明过了要证明的就是乘法幺半群同态

有乘法单位元 每一个环同态核都是理想 像都是子环

\end{lemma}

\begin{note}[关联图片]
以下图片已纳入本条整理，当前版本优先依据 Markdown 文字整理，图片细节可在审核时对照：

1. image.png\\
2. image 1.png\\

\end{note}

\begin{note}[来源路径]

近世代数/近世代数/无标题/引理2 3从理想可以从环同态中定义 1edd2efa3f4280168bf8f80d6cf1f8fc.md

\end{note}

\section*{引理2.7 引理2.8素理想与素数}

\begin{note}[来源上级主题]

环 (\%E7\%8E\%AF\%201d9d2efa3f4280f38e94de87f834c519.md)

\end{note}

\begin{lemma}[引理2.7 引理2.8素理想与素数]

5 |B 2.7

Gp RHER, abeZ, M

plab = > pla Xpl|b                        (2.125) 5

JERS aEDIT Gan HY, Bll

5 |B 2.8

BnK+hSR, WHE aDETZ, iF

n\textbackslash{}ab                                    (2.126)

nta                                    (2.127)

ntb                                    (2.128)

9

UEW) TERA BN. Hn Z-POR, RNAWM-PAEPADH n = ab, HP a,b \# +1. Mi nlab, TH

In| > lal, knta (HAB-PRERA-SR, WPAN AMALADT ST A-MA). Ant bd. eH,

BAM BEE HY x 5] E.

ARTE LTS SS |F, FA, CE, SHE RS AY. SEP LE, WR ER A BABE

A, XSLT DAC

Va,b € Z,(ab € pZ = > aé pZ Kb € pZ)                             (2.129)

FAG, FEAR ACHR, FR ae ER BE Ei COIR SCR FR EAE, Re

— NIC EAEIX A) REY

\end{lemma}

\begin{note}[图片 OCR 摘录，待人工复核]
以下内容来自源图片识别，便于审核时对照：

\textbf{图片 1}（image.png）
\begin{Verbatim}[fontsize=\small]
5 |B 2.7
Gp RHER, abeZ, M
plab = > pla Xpl|b                        (2.125) 5
JERS aEDIT Gan HY, Bll
5 |B 2.8
BnK+hSR, WHE aDETZ, iF
n\ab                                    (2.126)
nta                                    (2.127)
ntb                                    (2.128)
9
UEW) TERA BN. Hn Z-POR, RNAWM-PAEPADH n = ab, HP a,b # +1. Mi nlab, TH
In| > lal, knta (HAB-PRERA-SR, WPAN AMALADT ST A-MA). Ant bd. eH,
BAM BEE HY x 5] E.
ARTE LTS SS |F, FA, CE, SHE RS AY. SEP LE, WR ER A BABE
A, XSLT DAC
Va,b € Z,(ab € pZ = > aé pZ Kb € pZ)                             (2.129)
FAG, FEAR ACHR, FR ae ER BE Ei COIR SCR FR EAE, Re
— NIC EAEIX A) REY
\end{Verbatim}

\end{note}

\begin{note}[来源路径]

近世代数/近世代数/无标题/引理2 7 引理2 8素理想与素数 1d9d2efa3f42800ea8a7e5c78425a40c.md

\end{note}

\section*{循环子群}

\begin{note}[来源上级主题]

子群 (\%E5\%AD\%90\%E7\%BE\%A4\%201bdd2efa3f42801b985de1cf59e5ed5f.md)

\end{note}

\begin{note}[循环子群]

循环群的阶是父群阶的因数，阶也就是个数

定理3有限群的阶都是父群因子的阶

\end{note}

\begin{note}[关联图片]
以下图片已纳入本条整理，当前版本优先依据 Markdown 文字整理，图片细节可在审核时对照：

1. image.png\\

\end{note}

\begin{note}[来源路径]

近世代数/近世代数/无标题/循环子群 1bdd2efa3f428088afeae785ed4258af.md

\end{note}

\section*{循环置换}

\begin{note}[来源上级主题]

置换群 (\%E7\%BD\%AE\%E6\%8D\%A2\%E7\%BE\%A4\%201bad2efa3f4280228c17c29195a2c525.md)

\end{note}

\begin{note}[循环置换]

子级 项目: 定理2 n元置换可以写成循环置换的乘积 (\%E5\%AE\%9A\%E7\%90\%862\%20n\%E5\%85\%83\%E7\%BD\%AE\%E6\%8D\%A2\%E5\%8F\%AF\%E4\%BB\%A5\%E5\%86\%99\%E6\%88\%90\%E5\%BE\%AA\%E7\%8E\%AF\%E7\%BD\%AE\%E6\%8D\%A2\%E7\%9A\%84\%E4\%B9\%98\%E7\%A7\%AF\%201bad2efa3f428028af82e07c1d9e66d5.md)

3循环 第k个接回第一个

3个元都是循环置换

不变的映射不写

不是循环置换但是可以分组

可以写成乘积

两个循环不能有公共的东西加不相连的循环

\end{note}

\begin{note}[关联图片]
以下图片已纳入本条整理，当前版本优先依据 Markdown 文字整理，图片细节可在审核时对照：

1. image.png\\
2. image 1.png\\
3. image 2.png\\
4. image 3.png\\
5. image 4.png\\
6. image 5.png\\

\end{note}

\begin{note}[来源路径]

近世代数/近世代数/无标题/循环置换 1bad2efa3f4280109703cec9c3fb3605.md

\end{note}

\section*{循环群}

\begin{note}[来源上级主题]

循环群 (\%E5\%BE\%AA\%E7\%8E\%AF\%E7\%BE\%A4\%201c3d2efa3f42802a8213f9e0014506cb.md)

\end{note}

\begin{note}[循环群]

循环群的关键是找到固定元，由固定元可以生成所有的元素，m是加法做多少次也就是运算多少次

余k的可由余1的自+k次

\end{note}

\begin{note}[关联图片]
以下图片已纳入本条整理，当前版本优先依据 Markdown 文字整理，图片细节可在审核时对照：

1. image.png\\
2. IMG\_2886.jpeg\\

\end{note}

\begin{note}[来源路径]

近世代数/近世代数/无标题/循环群 1bcd2efa3f428005af56c60b88647173.md

\end{note}

\section*{循环群}

\begin{note}[来源上级主题]

群论 (\%E7\%BE\%A4\%E8\%AE\%BA\%201c3d2efa3f428064a77dd3ff465d3548.md)

\end{note}

\begin{note}[循环群]

子级 项目: 循环群的同构 (\%E5\%BE\%AA\%E7\%8E\%AF\%E7\%BE\%A4\%E7\%9A\%84\%E5\%90\%8C\%E6\%9E\%84\%201bcd2efa3f42806c9e26ebba7f88827c.md), 循环群 (\%E5\%BE\%AA\%E7\%8E\%AF\%E7\%BE\%A4\%201bcd2efa3f428005af56c60b88647173.md), 证明循环群相等 (\%E8\%AF\%81\%E6\%98\%8E\%E5\%BE\%AA\%E7\%8E\%AF\%E7\%BE\%A4\%E7\%9B\%B8\%E7\%AD\%89\%201bdd2efa3f42807cbf33cc12a1e46539.md), 非循环群Klein (\%E9\%9D\%9E\%E5\%BE\%AA\%E7\%8E\%AF\%E7\%BE\%A4Klein\%201cad2efa3f4280da96c7c92e94edac04.md)

\end{note}

\begin{note}[来源路径]

近世代数/近世代数/无标题/循环群 1c3d2efa3f42802a8213f9e0014506cb.md

\end{note}

\section*{循环群的同构}

\begin{note}[来源上级主题]

循环群 (\%E5\%BE\%AA\%E7\%8E\%AF\%E7\%BE\%A4\%201c3d2efa3f42802a8213f9e0014506cb.md)

\end{note}

\begin{note}[循环群的同构]

把生成元做映射

同构证明是映射证明同态+单+满

因为阶数是n所以+nq就会回到原点

还是证明一一对应

\end{note}

\begin{note}[关联图片]
以下图片已纳入本条整理，当前版本优先依据 Markdown 文字整理，图片细节可在审核时对照：

1. image.png\\
2. image 1.png\\
3. image 2.png\\
4. image 3.png\\
5. image 4.png\\

\end{note}

\begin{note}[来源路径]

近世代数/近世代数/无标题/循环群的同构 1bcd2efa3f42806c9e26ebba7f88827c.md

\end{note}

\section*{整数环中 素数生成的都是极大理想}

\begin{note}[来源上级主题]

极大理想的定义 (\%E6\%9E\%81\%E5\%A4\%A7\%E7\%90\%86\%E6\%83\%B3\%E7\%9A\%84\%E5\%AE\%9A\%E4\%B9\%89\%201efd2efa3f4280c6b178ee01ba7ba017.md)

\end{note}

\begin{note}[整数环中 素数生成的都是极大理想]

包含一个和p互素是这个理想包含由p生成的理想的必要条件也就是存在p生成不了的远，然后用互素条件就能把关键的1证明在这个包含p的理想里面，有1的理想就是整个环

\end{note}

\begin{note}[关联图片]
以下图片已纳入本条整理，当前版本优先依据 Markdown 文字整理，图片细节可在审核时对照：

1. image.png\\
2. image 1.png\\

\end{note}

\begin{note}[来源路径]

近世代数/近世代数/无标题/整数环中 素数生成的都是极大理想 1efd2efa3f428000b6b0d466f32f8a2f.md

\end{note}

\section*{整环定义}

\begin{note}[来源上级主题]

第三章 (\%E7\%AC\%AC\%E4\%B8\%89\%E7\%AB\%A0\%201ecd2efa3f4280c7a6b0ea0de416c09f.md)

\end{note}

\begin{definition}[整环定义]

没有零因子

定义2.18整环没有零因子是交换环

零因子就是本身不是0但是和不为零的相乘可以是0的元素

\end{definition}

\begin{note}[关联图片]
以下图片已纳入本条整理，当前版本优先依据 Markdown 文字整理，图片细节可在审核时对照：

1. image.png\\

\end{note}

\begin{note}[来源路径]

近世代数/近世代数/无标题/整环定义 1ecd2efa3f4280109d5bd2a376069406.md

\end{note}

\section*{无标题}

\begin{note}[来源上级主题]

环 (\%E7\%8E\%AF\%201d9d2efa3f4280f38e94de87f834c519.md)

\end{note}

\begin{note}[无标题]

待根据图片进一步人工校对。

\end{note}

\begin{note}[来源路径]

近世代数/近世代数/无标题/无标题 1edd2efa3f428093b94bfaf16f9b40fe.md

\end{note}

\section*{无标题}

\begin{note}[来源上级主题]

定理1除环只有两个理想0理想和单位理想 (\%E5\%AE\%9A\%E7\%90\%861\%E9\%99\%A4\%E7\%8E\%AF\%E5\%8F\%AA\%E6\%9C\%89\%E4\%B8\%A4\%E4\%B8\%AA\%E7\%90\%86\%E6\%83\%B30\%E7\%90\%86\%E6\%83\%B3\%E5\%92\%8C\%E5\%8D\%95\%E4\%BD\%8D\%E7\%90\%86\%E6\%83\%B3\%201edd2efa3f4280198c05df2fd987f2b8.md)

\end{note}

\begin{note}[无标题]

待根据图片进一步人工校对。

\end{note}

\begin{note}[来源路径]

近世代数/近世代数/无标题/无标题 1edd2efa3f42809db73ce9e5011f76a6.md

\end{note}

\section*{无标题}

\begin{note}[来源上级主题]

定义2.10理想 理想的符号也是▶️指向理想 (\%E5\%AE\%9A\%E4\%B9\%892\%2010\%E7\%90\%86\%E6\%83\%B3\%20\%E7\%90\%86\%E6\%83\%B3\%E7\%9A\%84\%E7\%AC\%A6\%E5\%8F\%B7\%E4\%B9\%9F\%E6\%98\%AF\%E2\%96\%B6\%EF\%B8\%8F\%E6\%8C\%87\%E5\%90\%91\%E7\%90\%86\%E6\%83\%B3\%201edd2efa3f4280cbb2a2d6978eb4e340.md)

\end{note}

\begin{note}[无标题]

待根据图片进一步人工校对。

\end{note}

\begin{note}[来源路径]

近世代数/近世代数/无标题/无标题 1edd2efa3f4280e7b9eedcae01414587.md

\end{note}

\section*{无标题}

\begin{note}[来源上级主题]

环 (\%E7\%8E\%AF\%201d9d2efa3f4280f38e94de87f834c519.md)

\end{note}

\begin{note}[无标题]

待根据图片进一步人工校对。

\end{note}

\begin{note}[来源路径]

近世代数/近世代数/无标题/无标题 1eed2efa3f428085bd11fd375cdd2524.md

\end{note}

\section*{无标题}

\begin{note}[来源上级主题]

环 (\%E7\%8E\%AF\%201d9d2efa3f4280f38e94de87f834c519.md)

\end{note}

\begin{note}[无标题]

待根据图片进一步人工校对。

\end{note}

\begin{note}[来源路径]

近世代数/近世代数/无标题/无标题 1efd2efa3f42800b8e33d5f514dca155.md

\end{note}

\section*{无标题}

\begin{note}[来源上级主题]

环 (\%E7\%8E\%AF\%201d9d2efa3f4280f38e94de87f834c519.md)

\end{note}

\begin{note}[无标题]

待根据图片进一步人工校对。

\end{note}

\begin{note}[来源路径]

近世代数/近世代数/无标题/无标题 1efd2efa3f42803eabaeca56db77592b.md

\end{note}

\section*{无标题}

\begin{note}[来源上级主题]

环 (\%E7\%8E\%AF\%201d9d2efa3f4280f38e94de87f834c519.md)

\end{note}

\begin{note}[无标题]

待根据图片进一步人工校对。

\end{note}

\begin{note}[来源路径]

近世代数/近世代数/无标题/无标题 1efd2efa3f42805597bcc24547e6283c.md

\end{note}

\section*{无标题}

\begin{note}[来源上级主题]

3.2整环 (3\%202\%E6\%95\%B4\%E7\%8E\%AF\%20209d2efa3f4280cda320c843bba252f6.md)

\end{note}

\begin{note}[无标题]

待根据图片进一步人工校对。

\end{note}

\begin{note}[来源路径]

近世代数/近世代数/无标题/无标题 209d2efa3f4280769b8efec54d64c727.md

\end{note}

\section*{无限阶的群和任意循环群同态}

\begin{note}[来源上级主题]

证明循环群相等 (\%E8\%AF\%81\%E6\%98\%8E\%E5\%BE\%AA\%E7\%8E\%AF\%E7\%BE\%A4\%E7\%9B\%B8\%E7\%AD\%89\%201bdd2efa3f42807cbf33cc12a1e46539.md)

\end{note}

\begin{note}[无限阶的群和任意循环群同态]

同态满就行不必单，多对一可以一对多不行

不是-m-n是a上面—

\end{note}

\begin{note}[关联图片]
以下图片已纳入本条整理，当前版本优先依据 Markdown 文字整理，图片细节可在审核时对照：

1. image.png\\

\end{note}

\begin{note}[来源路径]

近世代数/近世代数/无标题/无限阶的群和任意循环群同态 1bdd2efa3f4280a3a71ff10703363a56.md

\end{note}

\section*{极大理想与域}

\begin{note}[来源上级主题]

9极大理想 (9\%E6\%9E\%81\%E5\%A4\%A7\%E7\%90\%86\%E6\%83\%B3\%201efd2efa3f42802fa98fcbcc6a6ed8b2.md)

\end{note}

\begin{note}[极大理想与域]

极大理想的定义

有单位元的交换换

\end{note}

\begin{note}[图片 OCR 摘录，待人工复核]
以下内容来自源图片识别，便于审核时对照：

\textbf{图片 1}（image.png）
\begin{Verbatim}[fontsize=\small]
wesE: (BUERE— VA BALICH ACR,
GER, Ml
RIE ARO FER.
\end{Verbatim}

\textbf{图片 2}（image 1.png）
\begin{Verbatim}[fontsize=\small]
LRNB—B. ARAB TARGET. BT
—+THR, RATTUAA R W—TPRABMKEA—T+HR, (FER RTS
FEAR ILA BEARDS, ee A AR.
\end{Verbatim}

\end{note}

\begin{note}[来源路径]

近世代数/近世代数/无标题/极大理想与域 205d2efa3f4280fab07ae7a25e907d3b.md

\end{note}

\section*{极大理想引理1 对剩余类环充要}

\begin{note}[来源上级主题]

9极大理想 (9\%E6\%9E\%81\%E5\%A4\%A7\%E7\%90\%86\%E6\%83\%B3\%201efd2efa3f42802fa98fcbcc6a6ed8b2.md)

\end{note}

\begin{lemma}[极大理想引理1 对剩余类环充要]

子级 项目: 证明 (\%E8\%AF\%81\%E6\%98\%8E\%201efd2efa3f4280b59d4de3019bfebf87.md)

文本: 90

商环只有零理想和单位理想 当被除的理想是极大理想

这是一个标准结论：

> 在主理想整环中，一个主理想 (a) 是极大理想 ⟺ a 是不可约元素。

>

\end{lemma}

\begin{note}[关联图片]
以下图片已纳入本条整理，当前版本优先依据 Markdown 文字整理，图片细节可在审核时对照：

1. image.png\\

\end{note}

\begin{note}[来源路径]

近世代数/近世代数/无标题/极大理想引理1 对剩余类环充要 1efd2efa3f4280549cb7d1f5b63a5db3.md

\end{note}

\section*{极大理想的定义}

\begin{note}[来源上级主题]

9极大理想 (9\%E6\%9E\%81\%E5\%A4\%A7\%E7\%90\%86\%E6\%83\%B3\%201efd2efa3f42802fa98fcbcc6a6ed8b2.md)

\end{note}

\begin{definition}[极大理想的定义]

子级 项目: 整数环中 素数生成的都是极大理想 (\%E6\%95\%B4\%E6\%95\%B0\%E7\%8E\%AF\%E4\%B8\%AD\%20\%E7\%B4\%A0\%E6\%95\%B0\%E7\%94\%9F\%E6\%88\%90\%E7\%9A\%84\%E9\%83\%BD\%E6\%98\%AF\%E6\%9E\%81\%E5\%A4\%A7\%E7\%90\%86\%E6\%83\%B3\%201efd2efa3f428000b6b0d466f32f8a2f.md)

不等于父环，第二不被其他理想包含

两步走 I≠R

包含I的理想只有R

\end{definition}

\begin{note}[关联图片]
以下图片已纳入本条整理，当前版本优先依据 Markdown 文字整理，图片细节可在审核时对照：

1. image.png\\
2. image 1.png\\

\end{note}

\begin{note}[来源路径]

近世代数/近世代数/无标题/极大理想的定义 1efd2efa3f4280c6b178ee01ba7ba017.md

\end{note}

\section*{特征的定义 对加法所有非零元素共同的阶}

\begin{note}[来源上级主题]

第三章 (\%E7\%AC\%AC\%E4\%B8\%89\%E7\%AB\%A0\%201ecd2efa3f4280c7a6b0ea0de416c09f.md)

\end{note}

\begin{definition}[特征的定义 对加法所有非零元素共同的阶]

eM —PHREAFH R WARS 7c Hy AA Mal YY OO WA OE BEY BT AY AIR R

HY FE.

RH, —TRASATHH R WHEM REAR, BAER Bit M

TW 1) hake TAs RAS BE TE WR A PRE, AR AIK PIT LI BK AR

Xt. ARIE E— “MR BY Bik» AA Th A A 49 a a A EEE A. BL

HE FRAT GE 46 TE

\end{definition}

\begin{note}[图片 OCR 摘录，待人工复核]
以下内容来自源图片识别，便于审核时对照：

\textbf{图片 1}（image.png）
\begin{Verbatim}[fontsize=\small]
eM —PHREAFH R WARS 7c Hy AA Mal YY OO WA OE BEY BT AY AIR R
HY FE.

RH, —TRASATHH R WHEM REAR, BAER Bit M
TW 1) hake TAs RAS BE TE WR A PRE, AR AIK PIT LI BK AR
Xt. ARIE E— “MR BY Bik» AA Th A A 49 a a A EEE A. BL
HE FRAT GE 46 TE
\end{Verbatim}

\end{note}

\begin{note}[来源路径]

近世代数/近世代数/无标题/特征的定义 对加法所有非零元素共同的阶 1edd2efa3f4280ff9276db7f4d563977.md

\end{note}

\section*{环}

\begin{note}[环]

子级 项目: 环定义2.1 (\%E7\%8E\%AF\%E5\%AE\%9A\%E4\%B9\%892\%201\%201d8d2efa3f4280caad84e366546b232a.md), 定义2.2交换环 (\%E5\%AE\%9A\%E4\%B9\%892\%202\%E4\%BA\%A4\%E6\%8D\%A2\%E7\%8E\%AF\%201d8d2efa3f4280d4b681c6f9125a7bd9.md), 命题2.1负元 0元的运算 (\%E5\%91\%BD\%E9\%A2\%982\%201\%E8\%B4\%9F\%E5\%85\%83\%200\%E5\%85\%83\%E7\%9A\%84\%E8\%BF\%90\%E7\%AE\%97\%201d8d2efa3f4280408c76c66618c847b3.md), 命题2.2 零环 加法单位元与乘法单位元相等 (\%E5\%91\%BD\%E9\%A2\%982\%202\%20\%E9\%9B\%B6\%E7\%8E\%AF\%20\%E5\%8A\%A0\%E6\%B3\%95\%E5\%8D\%95\%E4\%BD\%8D\%E5\%85\%83\%E4\%B8\%8E\%E4\%B9\%98\%E6\%B3\%95\%E5\%8D\%95\%E4\%BD\%8D\%E5\%85\%83\%E7\%9B\%B8\%E7\%AD\%89\%201d8d2efa3f428072bc30db04dbb3356f.md), 定义2.3 有乘法可逆元素叫做单位除环的性质很好 (\%E5\%AE\%9A\%E4\%B9\%892\%203\%20\%E6\%9C\%89\%E4\%B9\%98\%E6\%B3\%95\%E5\%8F\%AF\%E9\%80\%86\%E5\%85\%83\%E7\%B4\%A0\%E5\%8F\%AB\%E5\%81\%9A\%E5\%8D\%95\%E4\%BD\%8D\%E9\%99\%A4\%E7\%8E\%AF\%E7\%9A\%84\%E6\%80\%A7\%E8\%B4\%A8\%E5\%BE\%88\%E5\%A5\%BD\%201d8d2efa3f4280ae9d4bc1dfa10ee023.md), 定义2.4除环非零元素都是单位 (\%E5\%AE\%9A\%E4\%B9\%892\%204\%E9\%99\%A4\%E7\%8E\%AF\%E9\%9D\%9E\%E9\%9B\%B6\%E5\%85\%83\%E7\%B4\%A0\%E9\%83\%BD\%E6\%98\%AF\%E5\%8D\%95\%E4\%BD\%8D\%201d8d2efa3f428059bd92f93a477bc7b1.md), 定义2.5域交换的除环——更好的环 (\%E5\%AE\%9A\%E4\%B9\%892\%205\%E5\%9F\%9F\%E4\%BA\%A4\%E6\%8D\%A2\%E7\%9A\%84\%E9\%99\%A4\%E7\%8E\%AF\%E2\%80\%94\%E2\%80\%94\%E6\%9B\%B4\%E5\%A5\%BD\%E7\%9A\%84\%E7\%8E\%AF\%201d8d2efa3f4280daabf6fdeddf87f1ff.md), 定义2.6子环 群是子群 幺半群是子幺半群 符号< (\%E5\%AE\%9A\%E4\%B9\%892\%206\%E5\%AD\%90\%E7\%8E\%AF\%20\%E7\%BE\%A4\%E6\%98\%AF\%E5\%AD\%90\%E7\%BE\%A4\%20\%E5\%B9\%BA\%E5\%8D\%8A\%E7\%BE\%A4\%E6\%98\%AF\%E5\%AD\%90\%E5\%B9\%BA\%E5\%8D\%8A\%E7\%BE\%A4\%20\%E7\%AC\%A6\%E5\%8F\%B7\%201d8d2efa3f42800a807fc311bee30d05.md), 定义2.7子集A生成环<A> (\%E5\%AE\%9A\%E4\%B9\%892\%207\%E5\%AD\%90\%E9\%9B\%86A\%E7\%94\%9F\%E6\%88\%90\%E7\%8E\%AF\%20A\%201d8d2efa3f4280f9b734fb70e49c8370.md), 定义2.8环的直积 (\%E5\%AE\%9A\%E4\%B9\%892\%208\%E7\%8E\%AF\%E7\%9A\%84\%E7\%9B\%B4\%E7\%A7\%AF\%201d9d2efa3f42800ea898d0b403b41d0d.md), 命题2.5环的直积还是环 (\%E5\%91\%BD\%E9\%A2\%982\%205\%E7\%8E\%AF\%E7\%9A\%84\%E7\%9B\%B4\%E7\%A7\%AF\%E8\%BF\%98\%E6\%98\%AF\%E7\%8E\%AF\%201d9d2efa3f42807da879eee76a3f0d7e.md), 定义2.9环同态 (\%E5\%AE\%9A\%E4\%B9\%892\%209\%E7\%8E\%AF\%E5\%90\%8C\%E6\%80\%81\%201d9d2efa3f428053a89bd4f291079c1e.md), 理想与正规子群 (\%E7\%90\%86\%E6\%83\%B3\%E4\%B8\%8E\%E6\%AD\%A3\%E8\%A7\%84\%E5\%AD\%90\%E7\%BE\%A4\%201edd2efa3f42800dad4cc01d4143bb5c.md), 定义2.10理想 理想的符号也是▶️指向理想 (\%E5\%AE\%9A\%E4\%B9\%892\%2010\%E7\%90\%86\%E6\%83\%B3\%20\%E7\%90\%86\%E6\%83\%B3\%E7\%9A\%84\%E7\%AC\%A6\%E5\%8F\%B7\%E4\%B9\%9F\%E6\%98\%AF\%E2\%96\%B6\%EF\%B8\%8F\%E6\%8C\%87\%E5\%90\%91\%E7\%90\%86\%E6\%83\%B3\%201edd2efa3f4280cbb2a2d6978eb4e340.md), 命题2.6交换环的理想左理想就是右理想 (\%E5\%91\%BD\%E9\%A2\%982\%206\%E4\%BA\%A4\%E6\%8D\%A2\%E7\%8E\%AF\%E7\%9A\%84\%E7\%90\%86\%E6\%83\%B3\%E5\%B7\%A6\%E7\%90\%86\%E6\%83\%B3\%E5\%B0\%B1\%E6\%98\%AF\%E5\%8F\%B3\%E7\%90\%86\%E6\%83\%B3\%201edd2efa3f4280d6b552d742f9f0113f.md), 命题2.7整环乘整数都是理想 (\%E5\%91\%BD\%E9\%A2\%982\%207\%E6\%95\%B4\%E7\%8E\%AF\%E4\%B9\%98\%E6\%95\%B4\%E6\%95\%B0\%E9\%83\%BD\%E6\%98\%AF\%E7\%90\%86\%E6\%83\%B3\%201edd2efa3f4280879de0f548a3792784.md), 引理2.3从理想可以从环同态中定义 (\%E5\%BC\%95\%E7\%90\%862\%203\%E4\%BB\%8E\%E7\%90\%86\%E6\%83\%B3\%E5\%8F\%AF\%E4\%BB\%A5\%E4\%BB\%8E\%E7\%8E\%AF\%E5\%90\%8C\%E6\%80\%81\%E4\%B8\%AD\%E5\%AE\%9A\%E4\%B9\%89\%201edd2efa3f4280168bf8f80d6cf1f8fc.md), 命题2.8 R-R‘环同态 核都是理想 像都是子环 (\%E5\%91\%BD\%E9\%A2\%982\%208\%20R-R\%E2\%80\%98\%E7\%8E\%AF\%E5\%90\%8C\%E6\%80\%81\%20\%E6\%A0\%B8\%E9\%83\%BD\%E6\%98\%AF\%E7\%90\%86\%E6\%83\%B3\%20\%E5\%83\%8F\%E9\%83\%BD\%E6\%98\%AF\%E5\%AD\%90\%E7\%8E\%AF\%201edd2efa3f4280c7aab3cfeaab813a46.md), 定义2.11商环 (\%E5\%AE\%9A\%E4\%B9\%892\%2011\%E5\%95\%86\%E7\%8E\%AF\%201eed2efa3f4280d08291e336ed049bc4.md), 命题2.9商环是环 当且仅当I是理想 (\%E5\%91\%BD\%E9\%A2\%982\%209\%E5\%95\%86\%E7\%8E\%AF\%E6\%98\%AF\%E7\%8E\%AF\%20\%E5\%BD\%93\%E4\%B8\%94\%E4\%BB\%85\%E5\%BD\%93I\%E6\%98\%AF\%E7\%90\%86\%E6\%83\%B3\%201eed2efa3f4280c88023ce1caa85f335.md), 引理2.7 引理2.8素理想与素数 (\%E5\%BC\%95\%E7\%90\%862\%207\%20\%E5\%BC\%95\%E7\%90\%862\%208\%E7\%B4\%A0\%E7\%90\%86\%E6\%83\%B3\%E4\%B8\%8E\%E7\%B4\%A0\%E6\%95\%B0\%201d9d2efa3f42800ea8a7e5c78425a40c.md), 定义2.18整环没有零因子是交换环 (\%E5\%AE\%9A\%E4\%B9\%892\%2018\%E6\%95\%B4\%E7\%8E\%AF\%E6\%B2\%A1\%E6\%9C\%89\%E9\%9B\%B6\%E5\%9B\%A0\%E5\%AD\%90\%E6\%98\%AF\%E4\%BA\%A4\%E6\%8D\%A2\%E7\%8E\%AF\%201d8d2efa3f428099bbfed4ab53ac0f69.md), 目标 (\%E7\%9B\%AE\%E6\%A0\%87\%201edd2efa3f4280719ba4ca9df0dbc4b5.md), 无标题 (\%E6\%97\%A0\%E6\%A0\%87\%E9\%A2\%98\%201edd2efa3f428093b94bfaf16f9b40fe.md), 无标题 (\%E6\%97\%A0\%E6\%A0\%87\%E9\%A2\%98\%201eed2efa3f428085bd11fd375cdd2524.md), 无标题 (\%E6\%97\%A0\%E6\%A0\%87\%E9\%A2\%98\%201efd2efa3f42800b8e33d5f514dca155.md), 无标题 (\%E6\%97\%A0\%E6\%A0\%87\%E9\%A2\%98\%201efd2efa3f42803eabaeca56db77592b.md), 无标题 (\%E6\%97\%A0\%E6\%A0\%87\%E9\%A2\%98\%201efd2efa3f42805597bcc24547e6283c.md)

文本: 42（55）

\end{note}

\begin{note}[来源路径]

近世代数/近世代数/无标题/环 1d9d2efa3f4280f38e94de87f834c519.md

\end{note}

\section*{环中逆元}

\begin{note}[来源上级主题]

第三章 (\%E7\%AC\%AC\%E4\%B8\%89\%E7\%AB\%A0\%201ecd2efa3f4280c7a6b0ea0de416c09f.md)

\end{note}

\begin{note}[环中逆元]

EM —PR MACH —P5c 6 MY ioca WSR, BE

ba =ab=1.

—S7a RS ARA—TPBI. AA, Kia AMM AMS’, HA

bab’ =b(ab’) =b1=b

=(ba)b’ =1b’ =b’.

\end{note}

\begin{note}[图片 OCR 摘录，待人工复核]
以下内容来自源图片识别，便于审核时对照：

\textbf{图片 1}（image.png）
\begin{Verbatim}[fontsize=\small]
EM —PR MACH —P5c 6 MY ioca WSR, BE
ba =ab=1.
—S7a RS ARA—TPBI. AA, Kia AMM AMS’, HA
bab’ =b(ab’) =b1=b
=(ba)b’ =1b’ =b’.
\end{Verbatim}

\end{note}

\begin{note}[来源路径]

近世代数/近世代数/无标题/环中逆元 1ecd2efa3f428081aa09e3990df0dc13.md

\end{note}

\section*{环定义2.1}

\begin{note}[来源上级主题]

环 (\%E7\%8E\%AF\%201d9d2efa3f4280f38e94de87f834c519.md)

\end{note}

\begin{definition}[环定义2.1]

文本: 36

环两种运算，加法有交换群的所有性质，乘法有幺半群的所有性质

加法可以有交换率有结合律有逆元，乘法只有结合律不需要有逆元，有逆元叫单位除了0都有单位就是除环

环的乘法对加法有分配律

\end{definition}

\begin{note}[关联图片]
以下图片已纳入本条整理，当前版本优先依据 Markdown 文字整理，图片细节可在审核时对照：

1. image.png\\

\end{note}

\begin{note}[来源路径]

近世代数/近世代数/无标题/环定义2 1 1d8d2efa3f4280caad84e366546b232a.md

\end{note}

\section*{环的定义}

\begin{note}[来源上级主题]

第三章 (\%E7\%AC\%AC\%E4\%B8\%89\%E7\%AB\%A0\%201ecd2efa3f4280c7a6b0ea0de416c09f.md)

\end{note}

\begin{definition}[环的定义]

EM —ThEAR UYBR—-TM, 248

I. R\#—P HF, R—-DBL, R WF—S AY BM SR BGS FR BE

MPR 5

I. R WFAA BABE HN RBS FR LE 5

Il. X*PRAGAACE:

a(bc)=(ab)c

WFR WERHH70a, b, c MMU:

« 64

\end{definition}

\begin{note}[图片 OCR 摘录，待人工复核]
以下内容来自源图片识别，便于审核时对照：

\textbf{图片 1}（image.png）
\begin{Verbatim}[fontsize=\small]
EM —ThEAR UYBR—-TM, 248
I. R#—P HF, R—-DBL, R WF—S AY BM SR BGS FR BE
MPR 5
I. R WFAA BABE HN RBS FR LE 5
Il. X*PRAGAACE:
a(bc)=(ab)c
WFR WERHH70a, b, c MMU:
« 64
\end{Verbatim}

\end{note}

\begin{note}[来源路径]

近世代数/近世代数/无标题/环的定义 1ecd2efa3f42801e8a2dfcc3ca0e824e.md

\end{note}

\section*{理想与正规子群}

\begin{note}[来源上级主题]

环 (\%E7\%8E\%AF\%201d9d2efa3f4280f38e94de87f834c519.md)

\end{note}

\begin{note}[理想与正规子群]

找到正规子群这样好的结构 为了让核一定是字环

为什么要定义理想

\end{note}

\begin{note}[关联图片]
以下图片已纳入本条整理，当前版本优先依据 Markdown 文字整理，图片细节可在审核时对照：

1. image.png\\
2. image 1.png\\
3. image 2.png\\

\end{note}

\begin{note}[来源路径]

近世代数/近世代数/无标题/理想与正规子群 1edd2efa3f42800dad4cc01d4143bb5c.md

\end{note}

\section*{理想前言}

\begin{note}[理想前言]

BOER, AME RTH EA FRE itscy), Alt ker(f) AEE H. VA, (BEAT HE EMT EL,

tetic HIE TRE TH, EEO HH) “TERT” AWE AIL PAGEANT (BARRE MF). Fe

VE PY “TERT RE” SLAP NY “SHAR”. kt —- MENA Za, SAR Aer Tie PO

MC. AREA, im(f) REELS CEH TI). WBA, RMA EE,

R/ker(f) \textasciitilde{} im(f)                                              (2.33)

\end{note}

\begin{note}[图片 OCR 摘录，待人工复核]
以下内容来自源图片识别，便于审核时对照：

\textbf{图片 1}（image.png）
\begin{Verbatim}[fontsize=\small]
BOER, AME RTH EA FRE itscy), Alt ker(f) AEE H. VA, (BEAT HE EMT EL,
tetic HIE TRE TH, EEO HH) “TERT” AWE AIL PAGEANT (BARRE MF). Fe
VE PY “TERT RE” SLAP NY “SHAR”. kt —- MENA Za, SAR Aer Tie PO
MC. AREA, im(f) REELS CEH TI). WBA, RMA EE,

R/ker(f) ~ im(f)                                              (2.33)
\end{Verbatim}

\textbf{图片 2}（image 1.png）
\begin{Verbatim}[fontsize=\small]
Vaa—=N                      iow
/           a           \
|                |               |         |
——                                                an                   = {1}
\     |      \     |
XK SF        YL AS
G                              G’                              G’
Pe] 2.1: BR ZEBCHLAR SEER.

SRA JLRS: iz UALR, IEW ALAS BEIM, ERS, EUR, EBL ERT
WHR.

BREE iE, RS “AERO” ABAD, WRC REE LEB APR RE LUE. 4 PETE
SALE EMOIREADE, PRATT HT DAB Py 8 Ze BU I BLE UU. CHL J RENEE, (ELE RAISE
RAW. RATA RVE, AAG MLA, ABE ES EMMA RE, AT TULARE LENO
RG, APRA PNY RA) TR. ESL, LE SERIAL, ma
BOMBA REVIT © TEES ROAR, PRET. SEBEL, PAHS, ALY
IMDM. FRADE AS , BUN TCR LRM, Bi UE.
\end{Verbatim}

\end{note}

\begin{note}[来源路径]

近世代数/近世代数/无标题/理想前言 1d9d2efa3f42801280b8cc5259b6c349.md

\end{note}

\section*{理想子环定义}

\begin{note}[来源上级主题]

7理想 (7\%E7\%90\%86\%E6\%83\%B3\%201edd2efa3f4280a88bfce32f4c292dfb.md)

\end{note}

\begin{definition}[理想子环定义]

任意a属于理想所以只要1也就是单位元在理想里面这个理想就是整个环 理想是原环所有元素和a相乘生成的元素组成的

\end{definition}

\begin{note}[关联图片]
以下图片已纳入本条整理，当前版本优先依据 Markdown 文字整理，图片细节可在审核时对照：

1. image.png\\
2. image 1.png\\
3. image 2.png\\

\end{note}

\begin{note}[来源路径]

近世代数/近世代数/无标题/理想子环定义 1edd2efa3f4280889eb4daa974cd0880.md

\end{note}

\section*{理想并一下还是理想}

\begin{note}[理想并一下还是理想]

EM 2.15

4 (R4,.) E-PHR, ALISIR, Vi

I+J=\{at+b:ael,be Js\}                       2.79)

rj 2.12

E(Ri4,°) Z—-PK, MLISR, NI+SEEABM, |

I+J<R                            (2.80) ‘

VED) WES RHROEEAW. RAN RAEAREY “RU” .

RU+J)=RI+RI C1+J                               (2.81)

(1+ JR=IR+JIRCI4+J                        (2.82)

I+J<R                              (2.83)

S30, FHAGAQA, FERRARA. IR SIOSER BE, RKRNAM HAA.

\end{note}

\begin{note}[图片 OCR 摘录，待人工复核]
以下内容来自源图片识别，便于审核时对照：

\textbf{图片 1}（image.png）
\begin{Verbatim}[fontsize=\small]
EM 2.15
4 (R4,.) E-PHR, ALISIR, Vi
I+J={at+b:ael,be Js}                       2.79)
rj 2.12
E(Ri4,°) Z—-PK, MLISR, NI+SEEABM, |
I+J<R                            (2.80) ‘
VED) WES RHROEEAW. RAN RAEAREY “RU” .
RU+J)=RI+RI C1+J                               (2.81)
(1+ JR=IR+JIRCI4+J                        (2.82)
I+J<R                              (2.83)
S30, FHAGAQA, FERRARA. IR SIOSER BE, RKRNAM HAA.
\end{Verbatim}

\end{note}

\begin{note}[来源路径]

近世代数/近世代数/无标题/理想并一下还是理想 209d2efa3f42808eac5df717a464eb32.md

\end{note}

\section*{生成理想符号（a1，a2）}

\begin{note}[来源上级主题]

7理想 (7\%E7\%90\%86\%E6\%83\%B3\%201edd2efa3f4280a88bfce32f4c292dfb.md)

\end{note}

\begin{note}[生成理想符号（a1，a2）]

EM MAY ai, a2, +, a, ERB. REMRRNAAS

(A159 Az5 ***5dm HEHM.

\end{note}

\begin{note}[图片 OCR 摘录，待人工复核]
以下内容来自源图片识别，便于审核时对照：

\textbf{图片 1}（image.png）
\begin{Verbatim}[fontsize=\small]
EM MAY ai, a2, +, a, ERB. REMRRNAAS
(A159 Az5 ***5dm HEHM.
\end{Verbatim}

\textbf{图片 2}（image 1.png）
\begin{Verbatim}[fontsize=\small]
RNTER R BERR m FIC, ars, Ons BARK m Hoc, RANE
—TBFU, HUASHAAUS mK
Sy ts, te +5 _ (5s; € (a;))
FEswA R Asc. BTU R W—SBR. KARA DIEM.
RNA UMERM NIC a fila’,
A=syts,teerts, (s;€(a;)),
a’=si tsp terts’, (5; €(a;)),
BBA, HF s;—si€(a,),
a—a’=(s,—s} + (sp—sh be tm 5m VEU
FFAMFR W—MER Ir Rit, HY rs;, sir€ (ai),
Ta=rsytrsztertrsm EN,
ar=sirtsorteerts,r EN,
WU DARE AE 21, doe ts Om Wee BE AR.
\end{Verbatim}

\end{note}

\begin{note}[来源路径]

近世代数/近世代数/无标题/生成理想符号（a1，a2） 1edd2efa3f4280ed86cbd75462c70e23.md

\end{note}

\section*{目标}

\begin{note}[来源上级主题]

环 (\%E7\%8E\%AF\%201d9d2efa3f4280f38e94de87f834c519.md)

\end{note}

\begin{note}[目标]

文本: 55

\end{note}

\begin{note}[来源路径]

近世代数/近世代数/无标题/目标 1edd2efa3f4280719ba4ca9df0dbc4b5.md

\end{note}

\section*{证明}

\begin{note}[来源上级主题]

命题2.8 R-R‘环同态 核都是理想 像都是子环 (\%E5\%91\%BD\%E9\%A2\%982\%208\%20R-R\%E2\%80\%98\%E7\%8E\%AF\%E5\%90\%8C\%E6\%80\%81\%20\%E6\%A0\%B8\%E9\%83\%BD\%E6\%98\%AF\%E7\%90\%86\%E6\%83\%B3\%20\%E5\%83\%8F\%E9\%83\%BD\%E6\%98\%AF\%E5\%AD\%90\%E7\%8E\%AF\%201edd2efa3f4280c7aab3cfeaab813a46.md)

\end{note}

\begin{proof}

ker(f)是加法的正规子群 这个理想有什么好处呢，就是因为理想这个东西可以把R中的东西全都乘一遍再纳入其中，所以理想中的东西肯定是有同一个性质的比如都是2的倍数什么的 ker（f）是R的理想。而同态可以帮助你拆括号，把和起来有的性质让拆开的也有

**理想本质上就是在环** R **中闭合某种“乘法吸收”结构，它迫使你一旦选了一个元素，就必须把它在整个环中“乘出来”的所有结果都纳入进去。 小于号就是子环**

\end{proof}

\begin{note}[关联图片]
以下图片已纳入本条整理，当前版本优先依据 Markdown 文字整理，图片细节可在审核时对照：

1. image.png\\

\end{note}

\begin{note}[来源路径]

近世代数/近世代数/无标题/证明 1edd2efa3f428054a04be17abb7ec9aa.md

\end{note}

\section*{证明}

\begin{note}[来源上级主题]

命题2.2 零环 加法单位元与乘法单位元相等 (\%E5\%91\%BD\%E9\%A2\%982\%202\%20\%E9\%9B\%B6\%E7\%8E\%AF\%20\%E5\%8A\%A0\%E6\%B3\%95\%E5\%8D\%95\%E4\%BD\%8D\%E5\%85\%83\%E4\%B8\%8E\%E4\%B9\%98\%E6\%B3\%95\%E5\%8D\%95\%E4\%BD\%8D\%E5\%85\%83\%E7\%9B\%B8\%E7\%AD\%89\%201d8d2efa3f428072bc30db04dbb3356f.md)

\end{note}

\begin{proof}

a=a*1然后把1换成0即可

引理1.1幺半群的逆元构成群

\end{proof}

\begin{note}[关联图片]
以下图片已纳入本条整理，当前版本优先依据 Markdown 文字整理，图片细节可在审核时对照：

1. image.png\\

\end{note}

\begin{note}[来源路径]

近世代数/近世代数/无标题/证明 1edd2efa3f428098a226ca9720be1aba.md

\end{note}

\section*{证明}

\begin{note}[来源上级主题]

命题2.9商环是环 当且仅当I是理想 (\%E5\%91\%BD\%E9\%A2\%982\%209\%E5\%95\%86\%E7\%8E\%AF\%E6\%98\%AF\%E7\%8E\%AF\%20\%E5\%BD\%93\%E4\%B8\%94\%E4\%BB\%85\%E5\%BD\%93I\%E6\%98\%AF\%E7\%90\%86\%E6\%83\%B3\%201eed2efa3f4280c88023ce1caa85f335.md)

\end{note}

\begin{proof}

证明乘法是良定义的就是代表元的选取不影响

商环商群对比

\end{proof}

\begin{note}[关联图片]
以下图片已纳入本条整理，当前版本优先依据 Markdown 文字整理，图片细节可在审核时对照：

1. image.png\\
2. image 1.png\\
3. image 2.png\\

\end{note}

\begin{note}[来源路径]

近世代数/近世代数/无标题/证明 1eed2efa3f42807fbefeddb9976e5a4c.md

\end{note}

\section*{证明}

\begin{note}[来源上级主题]

极大理想引理1 对剩余类环充要 (\%E6\%9E\%81\%E5\%A4\%A7\%E7\%90\%86\%E6\%83\%B3\%E5\%BC\%95\%E7\%90\%861\%20\%E5\%AF\%B9\%E5\%89\%A9\%E4\%BD\%99\%E7\%B1\%BB\%E7\%8E\%AF\%E5\%85\%85\%E8\%A6\%81\%201efd2efa3f4280549cb7d1f5b63a5db3.md)

\end{note}

\begin{proof}

iEAR RA o RHA R BR=R/A HAAWH.

RAVE EMAAR. BE UARA, DER WH, It

HB4O. WA, HS SHS, CoOL TMVMRRAVER HB, VERE

SUMARSFU, PRI

we, R AAS EE.

\end{proof}

\begin{note}[图片 OCR 摘录，待人工复核]
以下内容来自源图片识别，便于审核时对照：

\textbf{图片 1}（image.png）
\begin{Verbatim}[fontsize=\small]
iEAR RA o RHA R BR=R/A HAAWH.

RAVE EMAAR. BE UARA, DER WH, It
HB4O. WA, HS SHS, CoOL TMVMRRAVER HB, VERE
SUMARSFU, PRI
we, R AAS EE.
\end{Verbatim}

\end{note}

\begin{note}[来源路径]

近世代数/近世代数/无标题/证明 1efd2efa3f4280b59d4de3019bfebf87.md

\end{note}

\section*{证明}

\begin{note}[来源上级主题]

引理2 交换环只有零理想和单位理想R是一个域 (\%E5\%BC\%95\%E7\%90\%862\%20\%E4\%BA\%A4\%E6\%8D\%A2\%E7\%8E\%AF\%E5\%8F\%AA\%E6\%9C\%89\%E9\%9B\%B6\%E7\%90\%86\%E6\%83\%B3\%E5\%92\%8C\%E5\%8D\%95\%E4\%BD\%8D\%E7\%90\%86\%E6\%83\%B3R\%E6\%98\%AF\%E4\%B8\%80\%E4\%B8\%AA\%E5\%9F\%9F\%20202d2efa3f42801f8790d49e461b4508.md)

\end{note}

\begin{proof}

TERR RRNA R WER AO. a PAE RM EP (a) BRARSH,

FEA REQ=R. Am R Hi 1E (a). (A (a) HIMBA SM rare

ROWE, Pry

l1=a’a (a’ER).

RH, ROSAS TSH MBA, R —PM. IESE.

AL EP S| FERIA WB

\end{proof}

\begin{note}[图片 OCR 摘录，待人工复核]
以下内容来自源图片识别，便于审核时对照：

\textbf{图片 1}（image.png）
\begin{Verbatim}[fontsize=\small]
TERR RRNA R WER AO. a PAE RM EP (a) BRARSH,
FEA REQ=R. Am R Hi 1E (a). (A (a) HIMBA SM rare
ROWE, Pry

l1=a’a (a’ER).
RH, ROSAS TSH MBA, R —PM. IESE.

AL EP S| FERIA WB
\end{Verbatim}

\end{note}

\begin{note}[来源路径]

近世代数/近世代数/无标题/证明 202d2efa3f4280ba8435d49791bf53ee.md

\end{note}

\section*{证明循环群相等}

\begin{note}[来源上级主题]

循环群 (\%E5\%BE\%AA\%E7\%8E\%AF\%E7\%BE\%A4\%201c3d2efa3f42802a8213f9e0014506cb.md)

\end{note}

\begin{proof}

子级 项目: 无限阶的群和任意循环群同态 (\%E6\%97\%A0\%E9\%99\%90\%E9\%98\%B6\%E7\%9A\%84\%E7\%BE\%A4\%E5\%92\%8C\%E4\%BB\%BB\%E6\%84\%8F\%E5\%BE\%AA\%E7\%8E\%AF\%E7\%BE\%A4\%E5\%90\%8C\%E6\%80\%81\%201bdd2efa3f4280a3a71ff10703363a56.md)

子集+数量一样 保证阶数一样

也就是和12互素

\end{proof}

\begin{note}[关联图片]
以下图片已纳入本条整理，当前版本优先依据 Markdown 文字整理，图片细节可在审核时对照：

1. image.png\\

\end{note}

\begin{note}[来源路径]

近世代数/近世代数/无标题/证明循环群相等 1bdd2efa3f42807cbf33cc12a1e46539.md

\end{note}

\section*{除环}

\begin{note}[来源上级主题]

第三章 (\%E7\%AC\%AC\%E4\%B8\%89\%E7\%AB\%A0\%201ecd2efa3f4280c7a6b0ea0de416c09f.md)

\end{note}

\begin{note}[除环]

除环强调的是有逆元和单位元

定义2.4除环非零元素都是单位

除环因为有逆元所以没有零因子

证明一个环是除环也就是任意两个数乘起来只要两个都不是0结果就不能是0

\end{note}

\begin{note}[关联图片]
以下图片已纳入本条整理，当前版本优先依据 Markdown 文字整理，图片细节可在审核时对照：

1. image.png\\

\end{note}

\begin{note}[来源路径]

近世代数/近世代数/无标题/除环 1edd2efa3f428092b9f2f5fd34528e4a.md

\end{note}

\section*{非循环群Klein}

\begin{note}[来源上级主题]

循环群 (\%E5\%BE\%AA\%E7\%8E\%AF\%E7\%BE\%A4\%201c3d2efa3f42802a8213f9e0014506cb.md)

\end{note}

\begin{note}[非循环群Klein]

HIF 1: Klein PUR V:

EM:

Klein Pic te

V, = \{e, a, b, ab\},

we

a’? =b? =(ab)? =e, ab=ba.

Ate Vi RAGA:

- SV BEAR, UUGERTTR 9 EBEMBMB, BDV = (9)

B ord(g) = 4,

- (Be Vit, RT MT edt, HRTH a,b A ab AWN 2, Aut

WA EI — TCA MAT LAA 4,

» URGERTARBERET i, AER E EME.

\end{note}

\begin{note}[图片 OCR 摘录，待人工复核]
以下内容来自源图片识别，便于审核时对照：

\textbf{图片 1}（image.png）
\begin{Verbatim}[fontsize=\small]
HIF 1: Klein PUR V:

EM:

Klein Pic te

V, = {e, a, b, ab},

we

a’? =b? =(ab)? =e, ab=ba.

Ate Vi RAGA:

- SV BEAR, UUGERTTR 9 EBEMBMB, BDV = (9)
B ord(g) = 4,

- (Be Vit, RT MT edt, HRTH a,b A ab AWN 2, Aut
WA EI — TCA MAT LAA 4,

» URGERTARBERET i, AER E EME.
\end{Verbatim}

\end{note}

\begin{note}[来源路径]

近世代数/近世代数/无标题/非循环群Klein 1cad2efa3f4280da96c7c92e94edac04.md

\end{note}
